%!TEX program = xelatex
\documentclass[11pt,a4paper]{article}
\usepackage[margin=2.2cm]{geometry}
\usepackage{amsmath,amssymb,amsthm,mathtools}
\usepackage{bm}
\usepackage{enumitem}
\usepackage{booktabs}
\usepackage{array}
\usepackage{xcolor}
\usepackage{tikz}
\usetikzlibrary{arrows.meta,positioning}
\usepackage{hyperref}
\usepackage{fancyhdr}
\usepackage{xeCJK}

\hypersetup{
    colorlinks=true,
    linkcolor=blue!60!black,
    citecolor=blue!60!black,
    urlcolor=blue!60!black
}

\pagestyle{fancy}
\fancyhf{}
\rhead{\small Qiuzhen QE Computational \& Applied Math --- Fall 2025}
\lhead{\small Official Qualifying Exam}
\rfoot{\small Page \thepage}

\DeclareMathOperator*{\argmax}{arg\,max}
\DeclareMathOperator*{\argmin}{arg\,min}

\setlist[itemize]{topsep=3pt,itemsep=2pt,parsep=0pt}
\setlist[enumerate]{topsep=3pt,itemsep=2pt,parsep=0pt}

\newcommand{\examstatement}[1]{%
  \par\addvspace{0.85em}%
  \noindent{\bfseries #1}\par\nobreak\addvspace{0.28em}%
}

% Shared amsthm note/remark/theorem styles
% Shared math extras for qe-review.
% Requires already-loaded: amsmath, amssymb.
%
% Classic pitfall: AMS blackboard-bold fonts have no digit glyphs, so
% \mathbb{1} renders as overlapping garbage. Map the digit 1 to dsfont's
% \mathds{1}; leave \mathbb{R}, \mathbb{P}, etc. on AMS.
%
% Usage (path relative to the main .tex being compiled):
%   notes:  % Shared math extras for qe-review.
% Requires already-loaded: amsmath, amssymb.
%
% Classic pitfall: AMS blackboard-bold fonts have no digit glyphs, so
% \mathbb{1} renders as overlapping garbage. Map the digit 1 to dsfont's
% \mathds{1}; leave \mathbb{R}, \mathbb{P}, etc. on AMS.
%
% Usage (path relative to the main .tex being compiled):
%   notes:  % Shared math extras for qe-review.
% Requires already-loaded: amsmath, amssymb.
%
% Classic pitfall: AMS blackboard-bold fonts have no digit glyphs, so
% \mathbb{1} renders as overlapping garbage. Map the digit 1 to dsfont's
% \mathds{1}; leave \mathbb{R}, \mathbb{P}, etc. on AMS.
%
% Usage (path relative to the main .tex being compiled):
%   notes:  \input{../../common/qe-math-extras.tex}
%   exams:  \input{../../../common/qe-math-extras.tex}

\usepackage{dsfont}
\usepackage{pdftexcmds}

\makeatletter
\let\qe@amsmmathbb\mathbb
\renewcommand{\mathbb}[1]{%
  \ifnum\pdf@strcmp{\unexpanded{#1}}{1}=\z@
    \mathds{1}%
  \else
    \qe@amsmmathbb{#1}%
  \fi
}
\makeatother

% Preferred indicator macros (either style is safe after the remap above).
\newcommand{\bbone}{\mathds{1}}
\newcommand{\ind}{\mathbf{1}}

%   exams:  % Shared math extras for qe-review.
% Requires already-loaded: amsmath, amssymb.
%
% Classic pitfall: AMS blackboard-bold fonts have no digit glyphs, so
% \mathbb{1} renders as overlapping garbage. Map the digit 1 to dsfont's
% \mathds{1}; leave \mathbb{R}, \mathbb{P}, etc. on AMS.
%
% Usage (path relative to the main .tex being compiled):
%   notes:  \input{../../common/qe-math-extras.tex}
%   exams:  \input{../../../common/qe-math-extras.tex}

\usepackage{dsfont}
\usepackage{pdftexcmds}

\makeatletter
\let\qe@amsmmathbb\mathbb
\renewcommand{\mathbb}[1]{%
  \ifnum\pdf@strcmp{\unexpanded{#1}}{1}=\z@
    \mathds{1}%
  \else
    \qe@amsmmathbb{#1}%
  \fi
}
\makeatother

% Preferred indicator macros (either style is safe after the remap above).
\newcommand{\bbone}{\mathds{1}}
\newcommand{\ind}{\mathbf{1}}


\usepackage{dsfont}
\usepackage{pdftexcmds}

\makeatletter
\let\qe@amsmmathbb\mathbb
\renewcommand{\mathbb}[1]{%
  \ifnum\pdf@strcmp{\unexpanded{#1}}{1}=\z@
    \mathds{1}%
  \else
    \qe@amsmmathbb{#1}%
  \fi
}
\makeatother

% Preferred indicator macros (either style is safe after the remap above).
\newcommand{\bbone}{\mathds{1}}
\newcommand{\ind}{\mathbf{1}}

%   exams:  % Shared math extras for qe-review.
% Requires already-loaded: amsmath, amssymb.
%
% Classic pitfall: AMS blackboard-bold fonts have no digit glyphs, so
% \mathbb{1} renders as overlapping garbage. Map the digit 1 to dsfont's
% \mathds{1}; leave \mathbb{R}, \mathbb{P}, etc. on AMS.
%
% Usage (path relative to the main .tex being compiled):
%   notes:  % Shared math extras for qe-review.
% Requires already-loaded: amsmath, amssymb.
%
% Classic pitfall: AMS blackboard-bold fonts have no digit glyphs, so
% \mathbb{1} renders as overlapping garbage. Map the digit 1 to dsfont's
% \mathds{1}; leave \mathbb{R}, \mathbb{P}, etc. on AMS.
%
% Usage (path relative to the main .tex being compiled):
%   notes:  \input{../../common/qe-math-extras.tex}
%   exams:  \input{../../../common/qe-math-extras.tex}

\usepackage{dsfont}
\usepackage{pdftexcmds}

\makeatletter
\let\qe@amsmmathbb\mathbb
\renewcommand{\mathbb}[1]{%
  \ifnum\pdf@strcmp{\unexpanded{#1}}{1}=\z@
    \mathds{1}%
  \else
    \qe@amsmmathbb{#1}%
  \fi
}
\makeatother

% Preferred indicator macros (either style is safe after the remap above).
\newcommand{\bbone}{\mathds{1}}
\newcommand{\ind}{\mathbf{1}}

%   exams:  % Shared math extras for qe-review.
% Requires already-loaded: amsmath, amssymb.
%
% Classic pitfall: AMS blackboard-bold fonts have no digit glyphs, so
% \mathbb{1} renders as overlapping garbage. Map the digit 1 to dsfont's
% \mathds{1}; leave \mathbb{R}, \mathbb{P}, etc. on AMS.
%
% Usage (path relative to the main .tex being compiled):
%   notes:  \input{../../common/qe-math-extras.tex}
%   exams:  \input{../../../common/qe-math-extras.tex}

\usepackage{dsfont}
\usepackage{pdftexcmds}

\makeatletter
\let\qe@amsmmathbb\mathbb
\renewcommand{\mathbb}[1]{%
  \ifnum\pdf@strcmp{\unexpanded{#1}}{1}=\z@
    \mathds{1}%
  \else
    \qe@amsmmathbb{#1}%
  \fi
}
\makeatother

% Preferred indicator macros (either style is safe after the remap above).
\newcommand{\bbone}{\mathds{1}}
\newcommand{\ind}{\mathbf{1}}


\usepackage{dsfont}
\usepackage{pdftexcmds}

\makeatletter
\let\qe@amsmmathbb\mathbb
\renewcommand{\mathbb}[1]{%
  \ifnum\pdf@strcmp{\unexpanded{#1}}{1}=\z@
    \mathds{1}%
  \else
    \qe@amsmmathbb{#1}%
  \fi
}
\makeatother

% Preferred indicator macros (either style is safe after the remap above).
\newcommand{\bbone}{\mathds{1}}
\newcommand{\ind}{\mathbf{1}}


\usepackage{dsfont}
\usepackage{pdftexcmds}

\makeatletter
\let\qe@amsmmathbb\mathbb
\renewcommand{\mathbb}[1]{%
  \ifnum\pdf@strcmp{\unexpanded{#1}}{1}=\z@
    \mathds{1}%
  \else
    \qe@amsmmathbb{#1}%
  \fi
}
\makeatother

% Preferred indicator macros (either style is safe after the remap above).
\newcommand{\bbone}{\mathds{1}}
\newcommand{\ind}{\mathbf{1}}

% Shared amsthm environments for qe-review (minimal).
% Requires already-loaded: amsthm.
% Safe to \input after \usepackage{amsmath,amssymb,amsthm,...}.
%
% Shared numbering uses aliascnt so cleveref keys off the environment
% name (Lemma, Proposition, Exercise, ...) rather than the parent
% counter (Theorem, Definition, Remark). Without aliascnt, \cref{lem:...}
% prints ``Theorem 9.13'' instead of ``Lemma 9.13''.

\usepackage{aliascnt}

\theoremstyle{plain}
\newtheorem{theorem}{Theorem}[section]

\newaliascnt{lemma}{theorem}
\newtheorem{lemma}[lemma]{Lemma}
\aliascntresetthe{lemma}

\newaliascnt{proposition}{theorem}
\newtheorem{proposition}[proposition]{Proposition}
\aliascntresetthe{proposition}

\newaliascnt{corollary}{theorem}
\newtheorem{corollary}[corollary]{Corollary}
\aliascntresetthe{corollary}

\newaliascnt{claim}{theorem}
\newtheorem{claim}[claim]{Claim}
\aliascntresetthe{claim}

\newaliascnt{fact}{theorem}
\newtheorem{fact}[fact]{Fact}
\aliascntresetthe{fact}

\theoremstyle{definition}
\newtheorem{definition}{Definition}[section]

\newaliascnt{exercise}{definition}
\newtheorem{exercise}[exercise]{Exercise}
\aliascntresetthe{exercise}

\newaliascnt{assumption}{definition}
\newtheorem{assumption}[assumption]{Assumption}
\aliascntresetthe{assumption}

\newaliascnt{notation}{definition}
\newtheorem{notation}[notation]{Notation}
\aliascntresetthe{notation}

\newaliascnt{construction}{definition}
\newtheorem{construction}[construction]{Construction}
\aliascntresetthe{construction}

% Example: document-wide numbering (not per section / chapter)
\newtheorem{example}{Example}

\theoremstyle{remark}
\newtheorem{remark}{Remark}[section]
\newtheorem*{remark*}{Remark}

\newaliascnt{heuristic}{remark}
\newtheorem{heuristic}[heuristic]{Heuristic}
\aliascntresetthe{heuristic}

\newaliascnt{convention}{remark}
\newtheorem{convention}[convention]{Convention}
\aliascntresetthe{convention}

\newaliascnt{warning}{remark}
\newtheorem{warning}[warning]{Warning}
\aliascntresetthe{warning}

\newaliascnt{proofsketch}{remark}
\newtheorem{proofsketch}[proofsketch]{Proof sketch}
\aliascntresetthe{proofsketch}

\newaliascnt{checkpoint}{remark}
\newtheorem{checkpoint}[checkpoint]{Checkpoint}
\aliascntresetthe{checkpoint}

% Note: document-wide numbering (not per section)
\newtheorem{note}{Note}
\newtheorem*{note*}{Note}

\newaliascnt{examnote}{note}
\newtheorem{examnote}[examnote]{Exam note}
\aliascntresetthe{examnote}

\newaliascnt{study}{note}
\newtheorem{study}[study]{Study note}
\aliascntresetthe{study}

\newaliascnt{summary}{note}
\newtheorem{summary}[summary]{Summary}
\aliascntresetthe{summary}

\newtheorem{examproblem}{Problem}[section]
\newtheorem*{examproblem*}{Problem}

% Bilingual aliases
\newenvironment{dingli}{\begin{theorem}}{\end{theorem}}
\newenvironment{yinli}{\begin{lemma}}{\end{lemma}}
\newenvironment{mingti}{\begin{proposition}}{\end{proposition}}
\newenvironment{tuilun}{\begin{corollary}}{\end{corollary}}
\newenvironment{dingyi}{\begin{definition}}{\end{definition}}
\newenvironment{li}{\begin{example}}{\end{example}}
\newenvironment{zhu}{\begin{remark}}{\end{remark}}
\newenvironment{biji}{\begin{note}}{\end{note}}
\newenvironment{kaoshi}{\begin{examnote}}{\end{examnote}}

\renewcommand{\qedsymbol}{\ensuremath{\square}}

% Shared cleveref setup for qe-review.
% Load AFTER: hyperref, and AFTER all \newtheorem / amsthm environments.
% Shared theorem counters are aliased in qe-amsthm-envs.tex (aliascnt),
% otherwise \cref on a lemma/proposition prints ``Theorem''.
%
% Usage (path relative to the main .tex being compiled):
%   notes:  % Shared cleveref setup for qe-review.
% Load AFTER: hyperref, and AFTER all \newtheorem / amsthm environments.
% Shared theorem counters are aliased in qe-amsthm-envs.tex (aliascnt),
% otherwise \cref on a lemma/proposition prints ``Theorem''.
%
% Usage (path relative to the main .tex being compiled):
%   notes:  % Shared cleveref setup for qe-review.
% Load AFTER: hyperref, and AFTER all \newtheorem / amsthm environments.
% Shared theorem counters are aliased in qe-amsthm-envs.tex (aliascnt),
% otherwise \cref on a lemma/proposition prints ``Theorem''.
%
% Usage (path relative to the main .tex being compiled):
%   notes:  \input{../../common/qe-cleveref.tex}
%   exams:  \input{../../../common/qe-cleveref.tex}
%
% Prefer \cref / \Cref over \ref. Use \Cref at sentence start.
% Unnumbered sectioning (e.g. \paragraph with default secnumdepth)
% produces an empty cleveref type; map that to "Paragraph".

\usepackage[nameinlink,noabbrev]{cleveref}

% Fallback for unnumbered \paragraph / \subsubsection labels
\crefname{}{Paragraph}{Paragraphs}
\Crefname{}{Paragraph}{Paragraphs}

% Numbered theorem-like (plain)
\crefname{theorem}{Theorem}{Theorems}
\Crefname{theorem}{Theorem}{Theorems}
\crefname{lemma}{Lemma}{Lemmas}
\Crefname{lemma}{Lemma}{Lemmas}
\crefname{proposition}{Proposition}{Propositions}
\Crefname{proposition}{Proposition}{Propositions}
\crefname{corollary}{Corollary}{Corollaries}
\Crefname{corollary}{Corollary}{Corollaries}
\crefname{claim}{Claim}{Claims}
\Crefname{claim}{Claim}{Claims}
\crefname{fact}{Fact}{Facts}
\Crefname{fact}{Fact}{Facts}

% Definition-like
\crefname{definition}{Definition}{Definitions}
\Crefname{definition}{Definition}{Definitions}
\crefname{example}{Example}{Examples}
\Crefname{example}{Example}{Examples}
\crefname{exercise}{Exercise}{Exercises}
\Crefname{exercise}{Exercise}{Exercises}
\crefname{assumption}{Assumption}{Assumptions}
\Crefname{assumption}{Assumption}{Assumptions}
\crefname{notation}{Notation}{Notations}
\Crefname{notation}{Notation}{Notations}
\crefname{construction}{Construction}{Constructions}
\Crefname{construction}{Construction}{Constructions}

% Remark-like
\crefname{remark}{Remark}{Remarks}
\Crefname{remark}{Remark}{Remarks}
\crefname{heuristic}{Heuristic}{Heuristics}
\Crefname{heuristic}{Heuristic}{Heuristics}
\crefname{convention}{Convention}{Conventions}
\Crefname{convention}{Convention}{Conventions}
\crefname{warning}{Warning}{Warnings}
\Crefname{warning}{Warning}{Warnings}
\crefname{proofsketch}{Proof sketch}{Proof sketches}
\Crefname{proofsketch}{Proof sketch}{Proof sketches}
\crefname{checkpoint}{Checkpoint}{Checkpoints}
\Crefname{checkpoint}{Checkpoint}{Checkpoints}

% Document-wide notes / exam problems
\crefname{note}{Note}{Notes}
\Crefname{note}{Note}{Notes}
\crefname{examnote}{Exam note}{Exam notes}
\Crefname{examnote}{Exam note}{Exam notes}
\crefname{study}{Study note}{Study notes}
\Crefname{study}{Study note}{Study notes}
\crefname{summary}{Summary}{Summaries}
\Crefname{summary}{Summary}{Summaries}
\crefname{examproblem}{Problem}{Problems}
\Crefname{examproblem}{Problem}{Problems}

% Sectioning / floats (explicit, noabbrev-friendly)
\crefname{section}{Section}{Sections}
\Crefname{section}{Section}{Sections}
\crefname{subsection}{Subsection}{Subsections}
\Crefname{subsection}{Subsection}{Subsections}
\crefname{subsubsection}{Subsubsection}{Subsubsections}
\Crefname{subsubsection}{Subsubsection}{Subsubsections}
\crefname{paragraph}{Paragraph}{Paragraphs}
\Crefname{paragraph}{Paragraph}{Paragraphs}
\crefname{equation}{Equation}{Equations}
\Crefname{equation}{Equation}{Equations}
\crefname{figure}{Figure}{Figures}
\Crefname{figure}{Figure}{Figures}
\crefname{table}{Table}{Tables}
\Crefname{table}{Table}{Tables}
\crefname{enumi}{Item}{Items}
\Crefname{enumi}{Item}{Items}

%   exams:  % Shared cleveref setup for qe-review.
% Load AFTER: hyperref, and AFTER all \newtheorem / amsthm environments.
% Shared theorem counters are aliased in qe-amsthm-envs.tex (aliascnt),
% otherwise \cref on a lemma/proposition prints ``Theorem''.
%
% Usage (path relative to the main .tex being compiled):
%   notes:  \input{../../common/qe-cleveref.tex}
%   exams:  \input{../../../common/qe-cleveref.tex}
%
% Prefer \cref / \Cref over \ref. Use \Cref at sentence start.
% Unnumbered sectioning (e.g. \paragraph with default secnumdepth)
% produces an empty cleveref type; map that to "Paragraph".

\usepackage[nameinlink,noabbrev]{cleveref}

% Fallback for unnumbered \paragraph / \subsubsection labels
\crefname{}{Paragraph}{Paragraphs}
\Crefname{}{Paragraph}{Paragraphs}

% Numbered theorem-like (plain)
\crefname{theorem}{Theorem}{Theorems}
\Crefname{theorem}{Theorem}{Theorems}
\crefname{lemma}{Lemma}{Lemmas}
\Crefname{lemma}{Lemma}{Lemmas}
\crefname{proposition}{Proposition}{Propositions}
\Crefname{proposition}{Proposition}{Propositions}
\crefname{corollary}{Corollary}{Corollaries}
\Crefname{corollary}{Corollary}{Corollaries}
\crefname{claim}{Claim}{Claims}
\Crefname{claim}{Claim}{Claims}
\crefname{fact}{Fact}{Facts}
\Crefname{fact}{Fact}{Facts}

% Definition-like
\crefname{definition}{Definition}{Definitions}
\Crefname{definition}{Definition}{Definitions}
\crefname{example}{Example}{Examples}
\Crefname{example}{Example}{Examples}
\crefname{exercise}{Exercise}{Exercises}
\Crefname{exercise}{Exercise}{Exercises}
\crefname{assumption}{Assumption}{Assumptions}
\Crefname{assumption}{Assumption}{Assumptions}
\crefname{notation}{Notation}{Notations}
\Crefname{notation}{Notation}{Notations}
\crefname{construction}{Construction}{Constructions}
\Crefname{construction}{Construction}{Constructions}

% Remark-like
\crefname{remark}{Remark}{Remarks}
\Crefname{remark}{Remark}{Remarks}
\crefname{heuristic}{Heuristic}{Heuristics}
\Crefname{heuristic}{Heuristic}{Heuristics}
\crefname{convention}{Convention}{Conventions}
\Crefname{convention}{Convention}{Conventions}
\crefname{warning}{Warning}{Warnings}
\Crefname{warning}{Warning}{Warnings}
\crefname{proofsketch}{Proof sketch}{Proof sketches}
\Crefname{proofsketch}{Proof sketch}{Proof sketches}
\crefname{checkpoint}{Checkpoint}{Checkpoints}
\Crefname{checkpoint}{Checkpoint}{Checkpoints}

% Document-wide notes / exam problems
\crefname{note}{Note}{Notes}
\Crefname{note}{Note}{Notes}
\crefname{examnote}{Exam note}{Exam notes}
\Crefname{examnote}{Exam note}{Exam notes}
\crefname{study}{Study note}{Study notes}
\Crefname{study}{Study note}{Study notes}
\crefname{summary}{Summary}{Summaries}
\Crefname{summary}{Summary}{Summaries}
\crefname{examproblem}{Problem}{Problems}
\Crefname{examproblem}{Problem}{Problems}

% Sectioning / floats (explicit, noabbrev-friendly)
\crefname{section}{Section}{Sections}
\Crefname{section}{Section}{Sections}
\crefname{subsection}{Subsection}{Subsections}
\Crefname{subsection}{Subsection}{Subsections}
\crefname{subsubsection}{Subsubsection}{Subsubsections}
\Crefname{subsubsection}{Subsubsection}{Subsubsections}
\crefname{paragraph}{Paragraph}{Paragraphs}
\Crefname{paragraph}{Paragraph}{Paragraphs}
\crefname{equation}{Equation}{Equations}
\Crefname{equation}{Equation}{Equations}
\crefname{figure}{Figure}{Figures}
\Crefname{figure}{Figure}{Figures}
\crefname{table}{Table}{Tables}
\Crefname{table}{Table}{Tables}
\crefname{enumi}{Item}{Items}
\Crefname{enumi}{Item}{Items}

%
% Prefer \cref / \Cref over \ref. Use \Cref at sentence start.
% Unnumbered sectioning (e.g. \paragraph with default secnumdepth)
% produces an empty cleveref type; map that to "Paragraph".

\usepackage[nameinlink,noabbrev]{cleveref}

% Fallback for unnumbered \paragraph / \subsubsection labels
\crefname{}{Paragraph}{Paragraphs}
\Crefname{}{Paragraph}{Paragraphs}

% Numbered theorem-like (plain)
\crefname{theorem}{Theorem}{Theorems}
\Crefname{theorem}{Theorem}{Theorems}
\crefname{lemma}{Lemma}{Lemmas}
\Crefname{lemma}{Lemma}{Lemmas}
\crefname{proposition}{Proposition}{Propositions}
\Crefname{proposition}{Proposition}{Propositions}
\crefname{corollary}{Corollary}{Corollaries}
\Crefname{corollary}{Corollary}{Corollaries}
\crefname{claim}{Claim}{Claims}
\Crefname{claim}{Claim}{Claims}
\crefname{fact}{Fact}{Facts}
\Crefname{fact}{Fact}{Facts}

% Definition-like
\crefname{definition}{Definition}{Definitions}
\Crefname{definition}{Definition}{Definitions}
\crefname{example}{Example}{Examples}
\Crefname{example}{Example}{Examples}
\crefname{exercise}{Exercise}{Exercises}
\Crefname{exercise}{Exercise}{Exercises}
\crefname{assumption}{Assumption}{Assumptions}
\Crefname{assumption}{Assumption}{Assumptions}
\crefname{notation}{Notation}{Notations}
\Crefname{notation}{Notation}{Notations}
\crefname{construction}{Construction}{Constructions}
\Crefname{construction}{Construction}{Constructions}

% Remark-like
\crefname{remark}{Remark}{Remarks}
\Crefname{remark}{Remark}{Remarks}
\crefname{heuristic}{Heuristic}{Heuristics}
\Crefname{heuristic}{Heuristic}{Heuristics}
\crefname{convention}{Convention}{Conventions}
\Crefname{convention}{Convention}{Conventions}
\crefname{warning}{Warning}{Warnings}
\Crefname{warning}{Warning}{Warnings}
\crefname{proofsketch}{Proof sketch}{Proof sketches}
\Crefname{proofsketch}{Proof sketch}{Proof sketches}
\crefname{checkpoint}{Checkpoint}{Checkpoints}
\Crefname{checkpoint}{Checkpoint}{Checkpoints}

% Document-wide notes / exam problems
\crefname{note}{Note}{Notes}
\Crefname{note}{Note}{Notes}
\crefname{examnote}{Exam note}{Exam notes}
\Crefname{examnote}{Exam note}{Exam notes}
\crefname{study}{Study note}{Study notes}
\Crefname{study}{Study note}{Study notes}
\crefname{summary}{Summary}{Summaries}
\Crefname{summary}{Summary}{Summaries}
\crefname{examproblem}{Problem}{Problems}
\Crefname{examproblem}{Problem}{Problems}

% Sectioning / floats (explicit, noabbrev-friendly)
\crefname{section}{Section}{Sections}
\Crefname{section}{Section}{Sections}
\crefname{subsection}{Subsection}{Subsections}
\Crefname{subsection}{Subsection}{Subsections}
\crefname{subsubsection}{Subsubsection}{Subsubsections}
\Crefname{subsubsection}{Subsubsection}{Subsubsections}
\crefname{paragraph}{Paragraph}{Paragraphs}
\Crefname{paragraph}{Paragraph}{Paragraphs}
\crefname{equation}{Equation}{Equations}
\Crefname{equation}{Equation}{Equations}
\crefname{figure}{Figure}{Figures}
\Crefname{figure}{Figure}{Figures}
\crefname{table}{Table}{Tables}
\Crefname{table}{Table}{Tables}
\crefname{enumi}{Item}{Items}
\Crefname{enumi}{Item}{Items}

%   exams:  % Shared cleveref setup for qe-review.
% Load AFTER: hyperref, and AFTER all \newtheorem / amsthm environments.
% Shared theorem counters are aliased in qe-amsthm-envs.tex (aliascnt),
% otherwise \cref on a lemma/proposition prints ``Theorem''.
%
% Usage (path relative to the main .tex being compiled):
%   notes:  % Shared cleveref setup for qe-review.
% Load AFTER: hyperref, and AFTER all \newtheorem / amsthm environments.
% Shared theorem counters are aliased in qe-amsthm-envs.tex (aliascnt),
% otherwise \cref on a lemma/proposition prints ``Theorem''.
%
% Usage (path relative to the main .tex being compiled):
%   notes:  \input{../../common/qe-cleveref.tex}
%   exams:  \input{../../../common/qe-cleveref.tex}
%
% Prefer \cref / \Cref over \ref. Use \Cref at sentence start.
% Unnumbered sectioning (e.g. \paragraph with default secnumdepth)
% produces an empty cleveref type; map that to "Paragraph".

\usepackage[nameinlink,noabbrev]{cleveref}

% Fallback for unnumbered \paragraph / \subsubsection labels
\crefname{}{Paragraph}{Paragraphs}
\Crefname{}{Paragraph}{Paragraphs}

% Numbered theorem-like (plain)
\crefname{theorem}{Theorem}{Theorems}
\Crefname{theorem}{Theorem}{Theorems}
\crefname{lemma}{Lemma}{Lemmas}
\Crefname{lemma}{Lemma}{Lemmas}
\crefname{proposition}{Proposition}{Propositions}
\Crefname{proposition}{Proposition}{Propositions}
\crefname{corollary}{Corollary}{Corollaries}
\Crefname{corollary}{Corollary}{Corollaries}
\crefname{claim}{Claim}{Claims}
\Crefname{claim}{Claim}{Claims}
\crefname{fact}{Fact}{Facts}
\Crefname{fact}{Fact}{Facts}

% Definition-like
\crefname{definition}{Definition}{Definitions}
\Crefname{definition}{Definition}{Definitions}
\crefname{example}{Example}{Examples}
\Crefname{example}{Example}{Examples}
\crefname{exercise}{Exercise}{Exercises}
\Crefname{exercise}{Exercise}{Exercises}
\crefname{assumption}{Assumption}{Assumptions}
\Crefname{assumption}{Assumption}{Assumptions}
\crefname{notation}{Notation}{Notations}
\Crefname{notation}{Notation}{Notations}
\crefname{construction}{Construction}{Constructions}
\Crefname{construction}{Construction}{Constructions}

% Remark-like
\crefname{remark}{Remark}{Remarks}
\Crefname{remark}{Remark}{Remarks}
\crefname{heuristic}{Heuristic}{Heuristics}
\Crefname{heuristic}{Heuristic}{Heuristics}
\crefname{convention}{Convention}{Conventions}
\Crefname{convention}{Convention}{Conventions}
\crefname{warning}{Warning}{Warnings}
\Crefname{warning}{Warning}{Warnings}
\crefname{proofsketch}{Proof sketch}{Proof sketches}
\Crefname{proofsketch}{Proof sketch}{Proof sketches}
\crefname{checkpoint}{Checkpoint}{Checkpoints}
\Crefname{checkpoint}{Checkpoint}{Checkpoints}

% Document-wide notes / exam problems
\crefname{note}{Note}{Notes}
\Crefname{note}{Note}{Notes}
\crefname{examnote}{Exam note}{Exam notes}
\Crefname{examnote}{Exam note}{Exam notes}
\crefname{study}{Study note}{Study notes}
\Crefname{study}{Study note}{Study notes}
\crefname{summary}{Summary}{Summaries}
\Crefname{summary}{Summary}{Summaries}
\crefname{examproblem}{Problem}{Problems}
\Crefname{examproblem}{Problem}{Problems}

% Sectioning / floats (explicit, noabbrev-friendly)
\crefname{section}{Section}{Sections}
\Crefname{section}{Section}{Sections}
\crefname{subsection}{Subsection}{Subsections}
\Crefname{subsection}{Subsection}{Subsections}
\crefname{subsubsection}{Subsubsection}{Subsubsections}
\Crefname{subsubsection}{Subsubsection}{Subsubsections}
\crefname{paragraph}{Paragraph}{Paragraphs}
\Crefname{paragraph}{Paragraph}{Paragraphs}
\crefname{equation}{Equation}{Equations}
\Crefname{equation}{Equation}{Equations}
\crefname{figure}{Figure}{Figures}
\Crefname{figure}{Figure}{Figures}
\crefname{table}{Table}{Tables}
\Crefname{table}{Table}{Tables}
\crefname{enumi}{Item}{Items}
\Crefname{enumi}{Item}{Items}

%   exams:  % Shared cleveref setup for qe-review.
% Load AFTER: hyperref, and AFTER all \newtheorem / amsthm environments.
% Shared theorem counters are aliased in qe-amsthm-envs.tex (aliascnt),
% otherwise \cref on a lemma/proposition prints ``Theorem''.
%
% Usage (path relative to the main .tex being compiled):
%   notes:  \input{../../common/qe-cleveref.tex}
%   exams:  \input{../../../common/qe-cleveref.tex}
%
% Prefer \cref / \Cref over \ref. Use \Cref at sentence start.
% Unnumbered sectioning (e.g. \paragraph with default secnumdepth)
% produces an empty cleveref type; map that to "Paragraph".

\usepackage[nameinlink,noabbrev]{cleveref}

% Fallback for unnumbered \paragraph / \subsubsection labels
\crefname{}{Paragraph}{Paragraphs}
\Crefname{}{Paragraph}{Paragraphs}

% Numbered theorem-like (plain)
\crefname{theorem}{Theorem}{Theorems}
\Crefname{theorem}{Theorem}{Theorems}
\crefname{lemma}{Lemma}{Lemmas}
\Crefname{lemma}{Lemma}{Lemmas}
\crefname{proposition}{Proposition}{Propositions}
\Crefname{proposition}{Proposition}{Propositions}
\crefname{corollary}{Corollary}{Corollaries}
\Crefname{corollary}{Corollary}{Corollaries}
\crefname{claim}{Claim}{Claims}
\Crefname{claim}{Claim}{Claims}
\crefname{fact}{Fact}{Facts}
\Crefname{fact}{Fact}{Facts}

% Definition-like
\crefname{definition}{Definition}{Definitions}
\Crefname{definition}{Definition}{Definitions}
\crefname{example}{Example}{Examples}
\Crefname{example}{Example}{Examples}
\crefname{exercise}{Exercise}{Exercises}
\Crefname{exercise}{Exercise}{Exercises}
\crefname{assumption}{Assumption}{Assumptions}
\Crefname{assumption}{Assumption}{Assumptions}
\crefname{notation}{Notation}{Notations}
\Crefname{notation}{Notation}{Notations}
\crefname{construction}{Construction}{Constructions}
\Crefname{construction}{Construction}{Constructions}

% Remark-like
\crefname{remark}{Remark}{Remarks}
\Crefname{remark}{Remark}{Remarks}
\crefname{heuristic}{Heuristic}{Heuristics}
\Crefname{heuristic}{Heuristic}{Heuristics}
\crefname{convention}{Convention}{Conventions}
\Crefname{convention}{Convention}{Conventions}
\crefname{warning}{Warning}{Warnings}
\Crefname{warning}{Warning}{Warnings}
\crefname{proofsketch}{Proof sketch}{Proof sketches}
\Crefname{proofsketch}{Proof sketch}{Proof sketches}
\crefname{checkpoint}{Checkpoint}{Checkpoints}
\Crefname{checkpoint}{Checkpoint}{Checkpoints}

% Document-wide notes / exam problems
\crefname{note}{Note}{Notes}
\Crefname{note}{Note}{Notes}
\crefname{examnote}{Exam note}{Exam notes}
\Crefname{examnote}{Exam note}{Exam notes}
\crefname{study}{Study note}{Study notes}
\Crefname{study}{Study note}{Study notes}
\crefname{summary}{Summary}{Summaries}
\Crefname{summary}{Summary}{Summaries}
\crefname{examproblem}{Problem}{Problems}
\Crefname{examproblem}{Problem}{Problems}

% Sectioning / floats (explicit, noabbrev-friendly)
\crefname{section}{Section}{Sections}
\Crefname{section}{Section}{Sections}
\crefname{subsection}{Subsection}{Subsections}
\Crefname{subsection}{Subsection}{Subsections}
\crefname{subsubsection}{Subsubsection}{Subsubsections}
\Crefname{subsubsection}{Subsubsection}{Subsubsections}
\crefname{paragraph}{Paragraph}{Paragraphs}
\Crefname{paragraph}{Paragraph}{Paragraphs}
\crefname{equation}{Equation}{Equations}
\Crefname{equation}{Equation}{Equations}
\crefname{figure}{Figure}{Figures}
\Crefname{figure}{Figure}{Figures}
\crefname{table}{Table}{Tables}
\Crefname{table}{Table}{Tables}
\crefname{enumi}{Item}{Items}
\Crefname{enumi}{Item}{Items}

%
% Prefer \cref / \Cref over \ref. Use \Cref at sentence start.
% Unnumbered sectioning (e.g. \paragraph with default secnumdepth)
% produces an empty cleveref type; map that to "Paragraph".

\usepackage[nameinlink,noabbrev]{cleveref}

% Fallback for unnumbered \paragraph / \subsubsection labels
\crefname{}{Paragraph}{Paragraphs}
\Crefname{}{Paragraph}{Paragraphs}

% Numbered theorem-like (plain)
\crefname{theorem}{Theorem}{Theorems}
\Crefname{theorem}{Theorem}{Theorems}
\crefname{lemma}{Lemma}{Lemmas}
\Crefname{lemma}{Lemma}{Lemmas}
\crefname{proposition}{Proposition}{Propositions}
\Crefname{proposition}{Proposition}{Propositions}
\crefname{corollary}{Corollary}{Corollaries}
\Crefname{corollary}{Corollary}{Corollaries}
\crefname{claim}{Claim}{Claims}
\Crefname{claim}{Claim}{Claims}
\crefname{fact}{Fact}{Facts}
\Crefname{fact}{Fact}{Facts}

% Definition-like
\crefname{definition}{Definition}{Definitions}
\Crefname{definition}{Definition}{Definitions}
\crefname{example}{Example}{Examples}
\Crefname{example}{Example}{Examples}
\crefname{exercise}{Exercise}{Exercises}
\Crefname{exercise}{Exercise}{Exercises}
\crefname{assumption}{Assumption}{Assumptions}
\Crefname{assumption}{Assumption}{Assumptions}
\crefname{notation}{Notation}{Notations}
\Crefname{notation}{Notation}{Notations}
\crefname{construction}{Construction}{Constructions}
\Crefname{construction}{Construction}{Constructions}

% Remark-like
\crefname{remark}{Remark}{Remarks}
\Crefname{remark}{Remark}{Remarks}
\crefname{heuristic}{Heuristic}{Heuristics}
\Crefname{heuristic}{Heuristic}{Heuristics}
\crefname{convention}{Convention}{Conventions}
\Crefname{convention}{Convention}{Conventions}
\crefname{warning}{Warning}{Warnings}
\Crefname{warning}{Warning}{Warnings}
\crefname{proofsketch}{Proof sketch}{Proof sketches}
\Crefname{proofsketch}{Proof sketch}{Proof sketches}
\crefname{checkpoint}{Checkpoint}{Checkpoints}
\Crefname{checkpoint}{Checkpoint}{Checkpoints}

% Document-wide notes / exam problems
\crefname{note}{Note}{Notes}
\Crefname{note}{Note}{Notes}
\crefname{examnote}{Exam note}{Exam notes}
\Crefname{examnote}{Exam note}{Exam notes}
\crefname{study}{Study note}{Study notes}
\Crefname{study}{Study note}{Study notes}
\crefname{summary}{Summary}{Summaries}
\Crefname{summary}{Summary}{Summaries}
\crefname{examproblem}{Problem}{Problems}
\Crefname{examproblem}{Problem}{Problems}

% Sectioning / floats (explicit, noabbrev-friendly)
\crefname{section}{Section}{Sections}
\Crefname{section}{Section}{Sections}
\crefname{subsection}{Subsection}{Subsections}
\Crefname{subsection}{Subsection}{Subsections}
\crefname{subsubsection}{Subsubsection}{Subsubsections}
\Crefname{subsubsection}{Subsubsection}{Subsubsections}
\crefname{paragraph}{Paragraph}{Paragraphs}
\Crefname{paragraph}{Paragraph}{Paragraphs}
\crefname{equation}{Equation}{Equations}
\Crefname{equation}{Equation}{Equations}
\crefname{figure}{Figure}{Figures}
\Crefname{figure}{Figure}{Figures}
\crefname{table}{Table}{Tables}
\Crefname{table}{Table}{Tables}
\crefname{enumi}{Item}{Items}
\Crefname{enumi}{Item}{Items}

%
% Prefer \cref / \Cref over \ref. Use \Cref at sentence start.
% Unnumbered sectioning (e.g. \paragraph with default secnumdepth)
% produces an empty cleveref type; map that to "Paragraph".

\usepackage[nameinlink,noabbrev]{cleveref}

% Fallback for unnumbered \paragraph / \subsubsection labels
\crefname{}{Paragraph}{Paragraphs}
\Crefname{}{Paragraph}{Paragraphs}

% Numbered theorem-like (plain)
\crefname{theorem}{Theorem}{Theorems}
\Crefname{theorem}{Theorem}{Theorems}
\crefname{lemma}{Lemma}{Lemmas}
\Crefname{lemma}{Lemma}{Lemmas}
\crefname{proposition}{Proposition}{Propositions}
\Crefname{proposition}{Proposition}{Propositions}
\crefname{corollary}{Corollary}{Corollaries}
\Crefname{corollary}{Corollary}{Corollaries}
\crefname{claim}{Claim}{Claims}
\Crefname{claim}{Claim}{Claims}
\crefname{fact}{Fact}{Facts}
\Crefname{fact}{Fact}{Facts}

% Definition-like
\crefname{definition}{Definition}{Definitions}
\Crefname{definition}{Definition}{Definitions}
\crefname{example}{Example}{Examples}
\Crefname{example}{Example}{Examples}
\crefname{exercise}{Exercise}{Exercises}
\Crefname{exercise}{Exercise}{Exercises}
\crefname{assumption}{Assumption}{Assumptions}
\Crefname{assumption}{Assumption}{Assumptions}
\crefname{notation}{Notation}{Notations}
\Crefname{notation}{Notation}{Notations}
\crefname{construction}{Construction}{Constructions}
\Crefname{construction}{Construction}{Constructions}

% Remark-like
\crefname{remark}{Remark}{Remarks}
\Crefname{remark}{Remark}{Remarks}
\crefname{heuristic}{Heuristic}{Heuristics}
\Crefname{heuristic}{Heuristic}{Heuristics}
\crefname{convention}{Convention}{Conventions}
\Crefname{convention}{Convention}{Conventions}
\crefname{warning}{Warning}{Warnings}
\Crefname{warning}{Warning}{Warnings}
\crefname{proofsketch}{Proof sketch}{Proof sketches}
\Crefname{proofsketch}{Proof sketch}{Proof sketches}
\crefname{checkpoint}{Checkpoint}{Checkpoints}
\Crefname{checkpoint}{Checkpoint}{Checkpoints}

% Document-wide notes / exam problems
\crefname{note}{Note}{Notes}
\Crefname{note}{Note}{Notes}
\crefname{examnote}{Exam note}{Exam notes}
\Crefname{examnote}{Exam note}{Exam notes}
\crefname{study}{Study note}{Study notes}
\Crefname{study}{Study note}{Study notes}
\crefname{summary}{Summary}{Summaries}
\Crefname{summary}{Summary}{Summaries}
\crefname{examproblem}{Problem}{Problems}
\Crefname{examproblem}{Problem}{Problems}

% Sectioning / floats (explicit, noabbrev-friendly)
\crefname{section}{Section}{Sections}
\Crefname{section}{Section}{Sections}
\crefname{subsection}{Subsection}{Subsections}
\Crefname{subsection}{Subsection}{Subsections}
\crefname{subsubsection}{Subsubsection}{Subsubsections}
\Crefname{subsubsection}{Subsubsection}{Subsubsections}
\crefname{paragraph}{Paragraph}{Paragraphs}
\Crefname{paragraph}{Paragraph}{Paragraphs}
\crefname{equation}{Equation}{Equations}
\Crefname{equation}{Equation}{Equations}
\crefname{figure}{Figure}{Figures}
\Crefname{figure}{Figure}{Figures}
\crefname{table}{Table}{Tables}
\Crefname{table}{Table}{Tables}
\crefname{enumi}{Item}{Items}
\Crefname{enumi}{Item}{Items}

% PDF sidebar outline + printed TOC for QE exam transcripts.
% Load in the preamble AFTER hyperref and AFTER \newcommand{\examstatement}.
\hypersetup{
  unicode=true,
  bookmarks=true,
  bookmarksopen=true,
  bookmarksopenlevel=3,
  bookmarksnumbered=false,
  bookmarksdepth=3
}
\IfFileExists{bookmark.sty}{%
  \usepackage{bookmark}%
  \bookmarksetup{open,openlevel=3,numbered=false,depth=3}%
}{}
% Unnumbered in print, but still written to the TOC / PDF outline
% down to \subsubsection. Starred headings create neither.
\setcounter{secnumdepth}{-1}
\setcounter{tocdepth}{3}
\makeatletter
\@ifundefined{examstatement}{}{%
  \let\qe@origexamstatement\examstatement
  \renewcommand{\examstatement}[1]{%
    \par\phantomsection
    \addcontentsline{toc}{subsection}{#1}%
    \qe@origexamstatement{#1}%
  }%
}
\makeatother


\title{\textbf{QUALIFYING EXAM: APPLIED MATHEMATICS}\\[0.5em]\Large AUTUMN 2025}
\date{}
\author{}

\begin{document}

\maketitle
\vspace{-3em}

\tableofcontents

\section{Statement of the problems}

\examstatement{Problem 1. [10 points]}
Except LU-type factorizations, matrix inversion can be obtained by iterations. Given \(A \in \mathbb{R}^{n \times n}\), consider the sequence \(\{X_{k}\}\) generated by:

\[
X_{k+1} = 2X_{k} - X_{k} A X_{k}, \qquad k = 0, 1, 2, \dots.
\]

Prove:
\begin{enumerate}[label=(\alph*)]
    \item if \(\|I - A X_{0}\| < 1\), then \(A\) is nonsingular and \(X_{k} \to A^{-1}\) quadratically.
    \item there exists \(X_{0}\) such that \(\|I - A X_{0}\| < 1\), if and only if \(A\) is nonsingular.
\end{enumerate}

\examstatement{Problem 2. [20 points]}
For any continuous function \(f\), define \(\|f\|_{\infty} = \max_{x \in [-1,1]} |f(x)|\).

Define the \(n\)-th Chebyshev polynomial by: \(T_{0}(x) = 1\), \(T_{1}(x) = x\), \(T_{n+1}(x) = 2x T_{n}(x) - T_{n-1}(x)\).
\begin{enumerate}[label=(\alph*)]
    \item Show that
    \[
    T_{n}(x) =
    \begin{cases}
    \cos(n \arccos x), & |x| \le 1, \\
    \cosh(n \operatorname{arccosh} x), & |x| > 1.
    \end{cases}
    \]
    \item Prove that \(T_{n}\) solves
    \[
    \min_{\substack{\text{polynomial }p\\ \deg p\le n,\, p(1)=1}}\|p\|_{\infty}
    \]
    and that the minimal value is \(1\), which is attained at \(p = T_{n}\).
    \item For \(f \in C^{n+1}[-1,1]\), choose the interpolation nodes \(x_{0}, x_{1}, \ldots, x_{n}\) as the \(n+1\) zeros of \(T_{n+1}(x)\), and denote its degree-\(n\) interpolation polynomial by \(L_{n}\). Prove
    \[
    \|f - L_{n}\|_{\infty} \le \frac{\|f^{(n+1)}\|_{\infty}}{(n+1)!\, 2^{n}}.
    \]
    \item For \(f \in C^{2n+2}[-1,1]\), derive the Gauss quadrature formula of
    \[
    I(f) = \int_{-1}^{1} \frac{f(x)}{\sqrt{1-x^{2}}}\,\mathrm{d}x
    \]
    for \(n+1\) quadrature nodes, written by \(I_{n+1}(f)\), and prove
    \[
    \bigl|I(f) - I_{n+1}(f)\bigr| \le \frac{\pi \|f^{(2n+2)}\|_{\infty}}{(2n+2)!\, 2^{2n+1}}.
    \]
\end{enumerate}

\examstatement{Problem 3. [20 points]}
In many applications generalized eigenvalue problem is needed to be solved: finding \(x \in \mathbb{C}^{n}\) and \(\mu, \lambda \in \mathbb{C}\), s.t.\ \(x \neq 0\), \(|\mu|^{2} + |\lambda|^{2} = 1\), \(\mu A x = \lambda B x\) for given \(A, B \in \mathbb{C}^{n \times n}\). Here \(x\) and \((\lambda, \mu)\) are called the eigenvector and eigenvalue of the pair \((A, B)\) respectively.

If \(B\) is nonsingular, this problem is equivalent to the standard eigenvalue problem \(A B^{-1} y = \mu^{-1} \lambda y\).

For two unitary matrices \(Q, Z\), \(Q^{H} A Z = S\), \(Q^{H} B Z = T\) is called a UET performing on \((A, B)\), written as \(Q^{H}(A, B)Z = (S, T)\) for ease.
\begin{enumerate}[label=(\alph*)]
    \item Give an example that \((A, B)\) has more than \(n\) eigenvalues.
    \item Prove that there exists a UET such that \(Q^{H}(A, B)Z = (S, T)\), where \(S, T\) are upper triangular.
    \item Can you obtain the eigenvalues of \((A, B)\) from \((S, T)\)? How?
    \item Propose a method to construct a UET such that \(Q_{0}^{H}(A, B)Z_{0} = (A_{0}, B_{0})\), where \(A_{0}\) is upper Hessenberg, and \(B_{0}\) is upper triangular. For ease, you may assume \(n = 4\).
\end{enumerate}

\examstatement{Problem 4. [15 points]}
Construct the Du Fort-Frankel scheme for the diffusion equation in 2D \(u_{t} = u_{xx} + u_{yy}\) and discuss its consistency and stability.

\examstatement{Problem 5. [15 points]}
For the wave equation \(u_{tt} = u_{xx}\), analyze the stability of the scheme

\[
\frac{u_{j}^{n+1} - 2u_{j}^{n} + u_{j}^{n-1}}{\tau^{2}}
=
\frac{u_{j+1}^{n+1} - 2u_{j}^{n+1} + u_{j-1}^{n+1}}{4h^{2}}
+
\frac{u_{j+1}^{n} - 2u_{j}^{n} + u_{j-1}^{n}}{2h^{2}}
+
\frac{u_{j+1}^{n-1} - 2u_{j}^{n-1} + u_{j-1}^{n-1}}{4h^{2}}.
\]

\noindent\textbf{You can choose either 6 or 7 to answer. The points will be decided as \(\max(6, 7)\).}

\examstatement{Problem 6. [20 points]}
Consider the so-called Rayleigh oscillator

\[
y'' + y + \epsilon\Bigl[\tfrac13 (y')^{3} - y'\Bigr] = 0,
\]

with initial condition

\[
y(0) = 0,\qquad y'(0) = 2a,
\]

where \(y = y(t)\), \(\epsilon > 0\), \(a > 0\).
\begin{enumerate}[label=(\alph*)]
    \item (14 points) For a small \(\epsilon\), construct an approximation of the solution to the above problem which is valid for large \(t\).
    \item (4 points) What is the accuracy of this approxmation? State a conclusion and briefly explain it.
    \item (2 points) Plot the approximated orbits in the phase plane, i.e., \(y - y'\) plane for several different \(a\). And explain what you observe.
\end{enumerate}

\examstatement{Problem 7. [20 points]}
Consider the minimization

\[
\min_{x \in \mathbb{R}^{n}} \|Ax - b\|_{2} + \gamma \|x\|_{1},
\]

with \(A \in \mathbb{R}^{m \times n}\), \(b \in \mathbb{R}^{m}\), and \(\gamma > 0\).
\begin{enumerate}[label=(\alph*)]
    \item (8 points) Derive the Lagrange dual of the equivalent problem:
    \[
    \min_{x,y} \|y\|_{2} + \gamma \|x\|_{1}, \quad \text{s.t. } Ax - b = y,
    \]
    with \(x \in \mathbb{R}^{n}\) and \(y \in \mathbb{R}^{m}\).
    \item (8 points) Suppose \(Ax^{*} - b \neq 0\) where \(x^{*}\) is an optimal point. Define \(r = \dfrac{Ax^{*} - b}{\|Ax^{*} - b\|_{2}}\). Show that
    \[
    \|A^{\top} r\|_{\infty} \le \gamma, \qquad r^{\top} A x^{*} + \gamma \|x^{*}\|_{1} = 0.
    \]
    \item (4 points) Show that if \(\|a_{i}\|_{2} < \gamma\) where \(a_{i}\) is the \(i\)-th column of \(A\), then \(x^{*}_{i} = 0\).
\end{enumerate}

\noindent\rule{\linewidth}{0.4pt}

\section{Statement and solutions}

\subsection{Problem 1. [10 points]}
Except LU-type factorizations, matrix inversion can be obtained by iterations. Given \(A \in \mathbb{R}^{n \times n}\), consider the sequence \(\{X_{k}\}\) generated by:

\[
X_{k+1} = 2X_{k} - X_{k} A X_{k}, \qquad k = 0, 1, 2, \dots.
\]

Prove:
\begin{enumerate}[label=(\alph*)]
    \item if \(\|I - A X_{0}\| < 1\), then \(A\) is nonsingular and \(X_{k} \to A^{-1}\) quadratically.
    \item there exists \(X_{0}\) such that \(\|I - A X_{0}\| < 1\), if and only if \(A\) is nonsingular.
\end{enumerate}

\setcounter{section}{1}
\setcounter{theorem}{0}
\setcounter{definition}{0}
\setcounter{remark}{0}
\setcounter{note}{0}

\subsubsection{Preliminaries: LU-type inversion, the method the first sentence excludes}

The exam opens with a split of the numerical linear algebra syllabus, not with a definition of the recurrence. Direct methods invert \(A\) by factoring it into matrices that are cheap to invert, then reversing the product. Iterative methods start from a seed \(X_{0}\) and refine it. Problem~1 scores the second family (Newton--Schulz). The first family is the default algorithm, and the sentence is unintelligible until ``LU-type'' is a two-phase template rather than a slogan. The official CAM references for that template are Golub--Van~Loan \cite{golub96} and Trefethen--Bau \cite{trefethen97}.

\begin{center}
\begin{tikzpicture}[
  box/.style={draw, rounded corners=2pt, align=center, inner sep=5pt, font=\small, minimum height=1.05cm},
  arr/.style={-{Latex}, thick}
]
  \node[font=\small\bfseries, anchor=east] at (-0.15,1.55) {direct};
  \node[box] (A) at (1.35,1.55) {\(A\)};
  \node[box] (F) at (4.55,1.55) {easy factors};
  \node[box] (inv) at (7.85,1.55) {\(A^{-1}\)};
  \draw[arr] (A) -- node[above, font=\footnotesize] {Phase~I} (F);
  \draw[arr] (F) -- node[above, font=\footnotesize] {Phase~II} (inv);

  \node[font=\small\bfseries, anchor=east] at (-0.15,-0.15) {iterative};
  \node[box] (AX) at (1.35,-0.15) {\((A,X_{0})\)};
  \node[box] (S) at (4.55,-0.15) {Schulz};
  \node[box] (inv2) at (7.85,-0.15) {\(A^{-1}\)};
  \draw[arr] (AX) -- (S);
  \draw[arr] (S) -- (inv2);
\end{tikzpicture}
\end{center}

\paragraph{The two-phase template.}

Every factorization whose factors are triangular, orthogonal, diagonal, or a permutation inverts \(A\) by the same skeleton. The name ``LU-type'' in the exam sentence is this skeleton, not the single identity \(A=LU\).

\begin{construction}[LU-type inversion template]
\label{con:lu-type}
Let \(A\in\mathbb{R}^{n\times n}\) be nonsingular. An \emph{LU-type} inversion is any algorithm of the following two-phase form.

\emph{Phase~I (factor).} Compute

\[
A = F_{1} F_{2} \cdots F_{m}
\]

in which each \(F_{i}\) belongs to a class whose inverse is \(O(n^{2})\) (or cheaper): unit lower or upper triangular, orthogonal or unitary, diagonal, or a permutation.

\emph{Phase~II (apply).} Do not form the dense matrices \(F_{i}^{-1}\) and multiply them. For each column \(e_{j}\) of \(I\), solve

\[
F_{1} F_{2} \cdots F_{m}\, x_{j} = e_{j}
\]

by applying the easy inverse of \(F_{1}\), then of \(F_{2}\), and so on. Assemble

\[
A^{-1}
=
\bigl[ x_{1} \mid \cdots \mid x_{n} \bigr].
\]

Equivalently, in product language,

\[
A^{-1}
=
F_{m}^{-1} \cdots F_{2}^{-1} F_{1}^{-1},
\]

with each \(F_{i}^{-1}\) realized by substitution, transposition, reciprocal of a diagonal, or a row swap---not by a second factorization.
\end{construction}

The cheap inverses are completely explicit.

\begin{lemma}[Easy inverses]
\label{lem:easy-inv}
Let \(P\) be a permutation matrix, \(Q\) orthogonal, \(D=\operatorname{diag}(d_{1},\ldots,d_{n})\) with every \(d_{i}\neq 0\), \(L\) unit lower triangular, and \(U\) upper triangular with nonzero diagonal. Then \(P^{-1}=P^{T}\), \(Q^{-1}=Q^{T}\), and \(D^{-1}=\operatorname{diag}(d_{1}^{-1},\ldots,d_{n}^{-1})\). The triangular inverses are the substitution maps: given \(b\in\mathbb{R}^{n}\), the unique \(y\) with \(Ly=b\) is

\[
y_{i}
=
b_{i}
-
\sum_{j=1}^{i-1} \ell_{ij} y_{j},
\qquad
i=1,\ldots,n
\]

(forward substitution; no division, because \(\ell_{ii}=1\)), and the unique \(x\) with \(Ux=b\) is

\[
x_{i}
=
\frac{1}{u_{ii}}
\Biggl(
b_{i}
-
\sum_{j=i+1}^{n} u_{ij} x_{j}
\Biggr),
\qquad
i=n,\ldots,1
\]

(back substitution). Each map is \(n^{2}\) flops.
\end{lemma}

\begin{proof}
Permutation and orthogonal matrices satisfy \(P^{T}P=I\) and \(Q^{T}Q=I\) by definition. The diagonal claim is entrywise. For \(L\), the \(i\)-th equation of \(Ly=b\) is \(y_{i}+\sum_{j<i}\ell_{ij}y_{j}=b_{i}\), which rearranges to the displayed recurrence; well-definedness and uniqueness are immediate by induction on \(i\). For \(U\), start from row \(n\) and induct downwards; each division is legal because \(u_{ii}\neq 0\). Counting: the \(i\)-th forward step does \(i-1\) fused multiply-adds, so \(\sum_{i=1}^{n}(i-1)=\frac12 n(n-1)\) multiplications and as many additions, i.e.\ \(n^{2}\) flops in the usual model. Back substitution is the same count plus \(n\) divisions.
\end{proof}

\paragraph{LU as the prototype.}

Gaussian elimination with partial pivoting is the default Phase~I for a general square matrix. Every nonsingular \(A\) admits it \cite{trefethen97}.

\begin{definition}[GEPP / \(PA=LU\)]
\label{def:gepp}
A \emph{PLU factorization} of \(A\in\mathbb{R}^{n\times n}\) is a permutation matrix \(P\), a unit lower triangular matrix \(L\), and an upper triangular matrix \(U\) with

\[
PA = LU.
\]

Gaussian elimination with partial pivoting (GEPP) produces one. At step \(k=1,\ldots,n-1\): among rows \(k,\ldots,n\) of the current \(A\), swap the row whose pivot-column entry has largest absolute value into position \(k\); record the swap in \(P\); then, for each \(i>k\), store the multiplier \(\ell_{ik}=a_{ik}/a_{kk}\) in \(L\) and replace row \(i\) by row \(i-\ell_{ik}\cdot(\text{row }k)\). The remaining upper triangle is \(U\). The algorithm succeeds with every \(u_{kk}\neq 0\) if and only if \(A\) is nonsingular. Without pivoting, \(A=LU\) exists if and only if every leading principal minor is nonzero; the standard obstruction is

\[
\begin{pmatrix}
0 & 1 \\
1 & 0
\end{pmatrix}.
\]
\end{definition}

\begin{construction}[Inversion via GEPP]
\label{con:gepp-inv}
Input: nonsingular \(A\in\mathbb{R}^{n\times n}\). Output: \(X=A^{-1}\).

\begin{enumerate}[label=\textup{(\arabic*)}]
    \item \emph{Phase~I.} Compute \(PA=LU\) by \cref{def:gepp}.
    \item \emph{Phase~II.} For \(j=1,\ldots,n\), with \(e_{j}\) the \(j\)-th column of \(I\):
    \begin{enumerate}[label=\textup{(\alph*)}]
        \item permute the right-hand side: \(b \coloneqq P e_{j}\);
        \item forward substitution: solve \(L y = b\) by \cref{lem:easy-inv};
        \item back substitution: solve \(U x_{j} = y\) by \cref{lem:easy-inv}.
    \end{enumerate}
    \item Assemble \(X=\bigl[x_{1}\mid\cdots\mid x_{n}\bigr]\).
\end{enumerate}

In product language this is \(A^{-1}=U^{-1}L^{-1}P\). The \(n\) right-hand sides are the columns of \(P\), not of \(I\), because the factorization is of \(PA\). Equivalently one solves the matrix equation \(LUX=P\) by a single blocked forward substitution \(LY=P\) followed by a single blocked back substitution \(UX=Y\).
\end{construction}

The arithmetic is \(\frac23 n^{3}\) flops for Phase~I plus \(n\cdot n^{2}=n^{3}\) for the \(n\) triangular pairs, hence \(\frac53 n^{3}\) flops to form the inverse. Solving one linear system \(Ax=b\) after the same factorization is only \(O(n^{2})\) more. Forming \(A^{-1}\) to compute \(A^{-1}b\) is therefore both more expensive and, in floating point, typically less accurate than \cref{con:gepp-inv} stopped after a single pair of substitutions \cite{golub96,trefethen97}.

\begin{example}[\(2\times 2\), no pivoting required]
\label{ex:lu-2x2}
Take

\[
A
=
\begin{pmatrix}
2 & 1 \\
4 & 3
\end{pmatrix}.
\]

The \((2,1)\) multiplier is \(\ell_{21}=4/2=2\), so

\[
L
=
\begin{pmatrix}
1 & 0 \\
2 & 1
\end{pmatrix},
\qquad
U
=
\begin{pmatrix}
2 & 1 \\
0 & 1
\end{pmatrix},
\qquad
P=I.
\]

Column \(j=1\), right-hand side \(e_{1}=(1,0)^{T}\). Forward: \(y_{1}=1\), \(y_{2}=0-2\cdot 1=-2\). Back: \(x_{2}=-2\), then \(2x_{1}+1\cdot(-2)=1\), so \(x_{1}=\frac32\). Column \(j=2\), right-hand side \(e_{2}=(0,1)^{T}\). Forward: \(y_{1}=0\), \(y_{2}=1\). Back: \(x_{2}=1\), then \(2x_{1}+1=0\), so \(x_{1}=-\frac12\). Hence

\[
A^{-1}
=
\begin{pmatrix}
\tfrac32 & -\tfrac12 \\
-2 & 1
\end{pmatrix}.
\]

The same two substitution pairs, with \(P\neq I\), are the whole of \cref{con:gepp-inv} in every dimension.
\end{example}

\paragraph{The rest of the family.}

The other named factorizations are the same template with a different list of easy factors. Gauss--Jordan---row-reducing \(\bigl[A\mid I\bigr]\) to \(\bigl[I\mid A^{-1}\bigr]\)---is the same elimination family: still direct, still not iterative, and not more stable than \cref{con:gepp-inv}.

{\footnotesize
\begin{center}
\begin{tabular}{@{}>{\raggedright\arraybackslash}p{2.55cm}
                   >{\raggedright\arraybackslash}p{2.85cm}
                   >{\raggedright\arraybackslash}p{3.15cm}
                   >{\raggedright\arraybackslash}p{4.15cm}@{}}
\toprule
Phase~I & Hypothesis & Inverse formula & Phase~II \\
\midrule
\(PA=LU\) (GEPP) & \(A\) nonsingular & \(U^{-1}L^{-1}P\) & permute; \(n\) forward/back solves \\
\(A=R^{T}R\) (Cholesky) & \(A\) SPD & \(R^{-1}R^{-T}\) & two triangular solves per column \\
\(A=QR\) (Householder) & \(A\) nonsingular & \(R^{-1}Q^{T}\) & multiply by \(Q^{T}\); one back-sub \\
\(A=U\Sigma V^{T}\) (SVD) & \(A\) nonsingular & \(V\Sigma^{-1}U^{T}\) & reciprocal of singular values; two products \\
\bottomrule
\end{tabular}
\end{center}
}

Householder QR always exists. Cholesky exists if and only if \(A\) is symmetric positive definite; it is the SPD special case of LU, with \(P=I\) and \(L=R^{T}\) (up to the conventional unit-diagonal scaling of \(L\)). The SVD always exists, and is the stable way to invert a matrix that is close to singular: one discards or regularizes tiny \(\sigma_{i}\) instead of dividing by them. In all four rows, the inverse of an orthogonal factor is a transpose, the inverse of a triangular factor is \cref{lem:easy-inv}, and the inverse of a diagonal factor is entrywise reciprocal. That is the entire content of ``LU-type''.

\begin{remark}[What one actually computes]
\label{rem:dont-invert}
If the task is a linear system \(Ax=b\), stop after Phase~I plus one application of Phase~II to \(b\). The object \(A^{-1}\) is needed only when many right-hand sides are not known in advance, or when the inverse itself is the output (condition-number estimates, bilinear forms, the present exam). The 2026 Spring Kronecker / Gram-matrix warning is the same principle: factor \(A\), do not form \((A^{T}A)^{-1}\) by a second inversion \cite{trefethen97}.
\end{remark}

\paragraph{Where the iteration sits.}

The recurrence in the exam is the other machine. It never factors \(A\). It uses \(A\) only through the products \(AX_{k}\) and \(X_{k}(AX_{k})\), and it is Newton applied to \(X\mapsto X^{-1}-A\) (equivalently, Schulz's residual-squaring map \(E_{k+1}=E_{k}^{2}\) for \(E_{k}=I-AX_{k}\)). That is why the opening sentence is an exception clause: given a seed with \(\|I-AX_{0}\|<1\), one can invert without LU. The marks are (a)--(b) of that iteration, not a request to execute \cref{con:gepp-inv} on the script.

\begin{examnote}[What of this belongs on a script]
\label{examnote:p1-lu}
\Cref{con:lu-type,con:gepp-inv} are orientation for the first sentence. A grader scoring Problem~1 is reading the residual identity \(I-AX_{k+1}=(I-AX_{k})^{2}\), the Neumann argument that \(\|I-AX_{0}\|<1\) forces \(A\) nonsingular, and the quadratic rate. Writing GEPP in full is not required. The template is the reason the sentence mentions LU at all: every named dense factorization inverts \(A\) by Phase~I then \(n\) easy solves, and Schulz is what remains when that route is set aside.
\end{examnote}


\subsection{Problem 2. [20 points]}
For any continuous function \(f\), define \(\|f\|_{\infty} = \max_{x \in [-1,1]} |f(x)|\).

Define the \(n\)-th Chebyshev polynomial by: \(T_{0}(x) = 1\), \(T_{1}(x) = x\), \(T_{n+1}(x) = 2x T_{n}(x) - T_{n-1}(x)\).
\begin{enumerate}[label=(\alph*)]
    \item Show that
    \[
    T_{n}(x) =
    \begin{cases}
    \cos(n \arccos x), & |x| \le 1, \\
    \cosh(n \operatorname{arccosh} x), & |x| > 1.
    \end{cases}
    \]
    \item Prove that \(T_{n}\) solves
    \[
    \min_{\substack{\text{polynomial }p\\ \deg p\le n,\, p(1)=1}}\|p\|_{\infty}
    \]
    and that the minimal value is \(1\), which is attained at \(p = T_{n}\).
    \item For \(f \in C^{n+1}[-1,1]\), choose the interpolation nodes \(x_{0}, x_{1}, \ldots, x_{n}\) as the \(n+1\) zeros of \(T_{n+1}(x)\), and denote its degree-\(n\) interpolation polynomial by \(L_{n}\). Prove
    \[
    \|f - L_{n}\|_{\infty} \le \frac{\|f^{(n+1)}\|_{\infty}}{(n+1)!\, 2^{n}}.
    \]
    \item For \(f \in C^{2n+2}[-1,1]\), derive the Gauss quadrature formula of
    \[
    I(f) = \int_{-1}^{1} \frac{f(x)}{\sqrt{1-x^{2}}}\,\mathrm{d}x
    \]
    for \(n+1\) quadrature nodes, written by \(I_{n+1}(f)\), and prove
    \[
    \bigl|I(f) - I_{n+1}(f)\bigr| \le \frac{\pi \|f^{(2n+2)}\|_{\infty}}{(2n+2)!\, 2^{2n+1}}.
    \]
\end{enumerate}

\setcounter{section}{2}
\setcounter{theorem}{0}
\setcounter{definition}{0}
\setcounter{remark}{0}
\setcounter{note}{0}

\subsubsection{Preliminaries: what a quadrature formula is, and why (d) uses the nodes of (c)}

Problem~2 is one object used four times. Parts~(a)--(b) identify the Chebyshev polynomial \(T_{n}\). Parts~(c)--(d) use the zeros of the next polynomial \(T_{n+1}\): once as interpolation nodes, once as quadrature nodes. The word the exam does not define is the second of those.

A quadrature formula is a finite recipe that replaces an integral by a weighted sum of samples of \(f\). The sample locations are the \emph{nodes}; the coefficients in front of the samples are the \emph{weights}. ``Gauss'' is a rule for choosing the nodes, not a different kind of sum. The 2026 Spring Problem~2 is the same package with Legendre polynomials and weight \(1\); this paper is the Chebyshev-weight twin \cite{conte00}.

\begin{center}
\begin{tikzpicture}[
  box/.style={draw, rounded corners=2pt, align=center, inner sep=4.5pt, font=\small},
  arr/.style={-{Latex}, thick}
]
  \node[box] (Tn) at (0,1.35) {\(T_{n}\)};
  \node[box] (a) at (3.15,2.15) {(a) closed form};
  \node[box] (b) at (3.15,0.55) {(b) minimax, \(p(1)=1\)};
  \node[box] (z) at (0,-1.15) {zeros of \(T_{n+1}\)};
  \node[box] (c) at (3.15,-0.55) {(c) interpolation nodes};
  \node[box] (d) at (3.15,-1.75) {(d) quadrature nodes};
  \draw[arr] (Tn) -- (a);
  \draw[arr] (Tn) -- (b);
  \draw[arr] (Tn) -- (z);
  \draw[arr] (z) -- (c);
  \draw[arr] (z) -- (d);
\end{tikzpicture}
\end{center}

\paragraph{The objects.}

Throughout, write \(\mathcal{P}_{k}\) for polynomials of degree at most \(k\), and write

\[
I(f)
\coloneqq
\int_{-1}^{1}
\frac{f(x)}{\sqrt{1-x^{2}}}\,
\mathrm{d}x
\]

for the weighted integral in (d). The weight \(w(x)=(1-x^{2})^{-1/2}\) is part of \(I\), not a factor one samples at the nodes.

\begin{definition}[Quadrature formula, nodes, weights]
\label{def:quadrature}
A \emph{quadrature formula} with \(n+1\) nodes for \(I\) is a linear functional of the form

\begin{equation}
\label{eq:quad-sum}
I_{n+1}(f)
\coloneqq
\sum_{j=0}^{n} w_{j}\, f(x_{j}).
\end{equation}

The points \(x_{0},\ldots,x_{n}\in[-1,1]\) are the \emph{quadrature nodes}; the scalars \(w_{0},\ldots,w_{n}\) are the \emph{quadrature weights}. The exam's phrase ``for \(n+1\) quadrature nodes, written by \(I_{n+1}(f)\)'' means: produce an explicit sum \eqref{eq:quad-sum}, i.e.\ name both the \(x_{j}\) and the \(w_{j}\). It does \emph{not} mean ``sample \(f(x)/\sqrt{1-x^{2}}\)''. The weight has already been integrated into the \(w_{j}\).
\end{definition}

The trapezoid rule, Simpson's rule, and the midpoint rule are quadrature formulae. Each is \eqref{eq:quad-sum} with a particular choice of nodes and weights. Gauss quadrature is the same shape, with a particular choice of nodes.

\begin{example}[One node, already Gauss]
\label{ex:gauss-cheb-one}
The unique node in \((-1,1)\) that can appear for \(n=0\) is the zero of \(T_{1}(x)=x\), namely \(x_{0}=0\). The only weight compatible with exactness on constants is \(w_{0}=I(1)=\pi\). Thus

\[
I_{1}(f)
=
\pi f(0).
\]

This is the statement ``\(I(f)\approx \pi f(0)\)'': one samples \(f\) at the origin, and the Chebyshev weight contributes the factor \(\pi=\int_{-1}^{1}(1-x^{2})^{-1/2}\,\mathrm{d}x\).
\end{example}

\paragraph{How the weights are determined, once the nodes are chosen.}

Fix any distinct nodes. There is a canonical way to produce weights so that \eqref{eq:quad-sum} integrates every polynomial of degree \(\le n\) exactly: integrate the interpolant.

\begin{definition}[Interpolatory quadrature]
\label{def:interpolatory}
Let \(x_{0},\ldots,x_{n}\) be distinct. Write \(\ell_{j}\) for the Lagrange basis

\[
\ell_{j}(x)
=
\prod_{\substack{0\le k\le n\\ k\neq j}}
\frac{x-x_{k}}{x_{j}-x_{k}},
\]

so that the unique \(L_{n}\in\mathcal{P}_{n}\) interpolating \(f\) at these nodes is \(L_{n}=\sum_{j} f(x_{j})\,\ell_{j}\). The associated \emph{interpolatory quadrature} is \eqref{eq:quad-sum} with

\[
w_{j}
\coloneqq
I(\ell_{j})
=
\int_{-1}^{1}
\frac{\ell_{j}(x)}{\sqrt{1-x^{2}}}\,
\mathrm{d}x.
\]

Equivalently, \(I_{n+1}(f)=I(L_{n})\): replace \(f\) by its interpolant, then integrate the interpolant exactly.
\end{definition}

\begin{lemma}[Exactness on \(\mathcal{P}_{n}\)]
\label{lem:interp-exact}
The interpolatory rule of \cref{def:interpolatory} satisfies \(I_{n+1}(p)=I(p)\) for every \(p\in\mathcal{P}_{n}\).
\end{lemma}

\begin{proof}
If \(\deg p\le n\), then \(L_{n}=p\), so \(I_{n+1}(p)=I(L_{n})=I(p)\).
\end{proof}

Part~(c) stops here: it asks for \(\|f-L_{n}\|_{\infty}\), not for \(I(L_{n})\). Part~(d) integrates that interpolant (in fact a Hermite interpolant of twice the degree, which is why the derivative order jumps). The nodes in (c) and (d) are the same points, used for two different linear maps.

\paragraph{What ``Gauss'' adds: a better choice of nodes.}

For generic nodes, \cref{lem:interp-exact} is sharp: some polynomial of degree \(n+1\) is missed. If the nodes are the zeros of the orthogonal polynomial of degree \(n+1\) for the inner product \(\langle f,g\rangle_{w}=I(fg)\), exactness doubles.

\begin{definition}[Gauss quadrature]
\label{def:gauss}
A quadrature \eqref{eq:quad-sum} with \(n+1\) nodes is \emph{Gauss} for the weight \(w\) if it is interpolatory in the sense of \cref{def:interpolatory} and the nodes are the zeros of a polynomial \(\pi_{n+1}\) of degree \(n+1\) that is orthogonal to \(\mathcal{P}_{n}\) in \(\langle\cdot,\cdot\rangle_{w}\). Equivalently: among all interpolatory rules with \(n+1\) nodes, Gauss is the unique one that integrates every element of \(\mathcal{P}_{2n+1}\) exactly.
\end{definition}

\begin{lemma}[Precision \(2n+1\)]
\label{lem:gauss-prec-cheb}
Let \(\pi_{n+1}\) be orthogonal to \(\mathcal{P}_{n}\) for \(\langle\cdot,\cdot\rangle_{w}\), with \(n+1\) distinct zeros \(x_{0},\ldots,x_{n}\) in \((-1,1)\). The interpolatory rule on those nodes satisfies \(I_{n+1}(p)=I(p)\) for every \(p\in\mathcal{P}_{2n+1}\).
\end{lemma}

\begin{proof}
Write \(p=q\pi_{n+1}+s\) with \(q,s\in\mathcal{P}_{n}\). Then \(I(q\pi_{n+1})=\langle q,\pi_{n+1}\rangle_{w}=0\), while \(p(x_{j})=s(x_{j})\), so \(I_{n+1}(p)=I_{n+1}(s)\). \Cref{lem:interp-exact} gives \(I_{n+1}(s)=I(s)\). Combining, \(I_{n+1}(p)=I(s)=I(p)\).
\end{proof}

That is the whole content of the adjective Gauss: spend the \(n+1\) free nodes so that the \(n+1\) interpolatory weights become exact on a space of dimension \(2n+2\). The error for a general \(f\in C^{2n+2}\) then involves \(f^{(2n+2)}\), which is why (d) assumes that regularity and why the factorial is \((2n+2)!\), not \((n+1)!\).

\paragraph{The Chebyshev case.}

After \(x=\cos\theta\), one has \(\mathrm{d}x/\sqrt{1-x^{2}}=-\mathrm{d}\theta\) and

\[
I(f)
=
\int_{0}^{\pi} f(\cos\theta)\,\mathrm{d}\theta.
\]

The Chebyshev polynomials \(T_{n}(\cos\theta)=\cos(n\theta)\) (the identity (a) asks to prove) are orthogonal on \([0,\pi]\) as ordinary cosines, hence orthogonal on \([-1,1]\) for the weight \((1-x^{2})^{-1/2}\). Their zeros are the interior points

\begin{equation}
\label{eq:cheb-nodes}
x_{j}
=
\cos\theta_{j},
\qquad
\theta_{j}
=
\frac{(2j+1)\pi}{2(n+1)},
\qquad
j=0,\ldots,n,
\end{equation}

which are exactly the interpolation nodes of (c). By \cref{def:gauss} they are the Gauss nodes of (d).

The angles \(\theta_{j}\) are the midpoints of \(n+1\) equal subintervals of \([0,\pi]\). The composite midpoint rule on those intervals is

\begin{equation}
\label{eq:gauss-cheb}
I_{n+1}(f)
=
\frac{\pi}{n+1}
\sum_{j=0}^{n} f(x_{j}),
\end{equation}

i.e.\ all weights equal \(w_{j}=\pi/(n+1)\). This is the Gauss quadrature formula (d) asks to derive: \(n+1\) nodes \eqref{eq:cheb-nodes}, and the constant weight \(\pi/(n+1)\). Exactness on \(\mathcal{P}_{2n+1}\) is \cref{lem:gauss-prec-cheb}, or equivalently exactness of the midpoint rule on trigonometric polynomials of degree \(\le 2n+1\) in \(\theta\).

\begin{example}[Two nodes]
\label{ex:gauss-cheb-two}
Here \(T_{2}(x)=2x^{2}-1\), with zeros \(x=\pm 1/\sqrt{2}\). Formula \eqref{eq:gauss-cheb} is

\[
I_{2}(f)
=
\frac{\pi}{2}
\Bigl(
f\bigl(\tfrac{1}{\sqrt{2}}\bigr)
+
f\bigl(-\tfrac{1}{\sqrt{2}}\bigr)
\Bigr).
\]

It is exact on \(\mathcal{P}_{3}\). It is not the trapezoid rule on \([-1,1]\) (that would sample the endpoints \(\pm 1\)).
\end{example}

\paragraph{The printed error constant.}

Let \(\omega(x)=\prod_{j=0}^{n}(x-x_{j})\). The leading coefficient of \(T_{n+1}\) is \(2^{n}\), so \(\omega=2^{-n}T_{n+1}\). If \(H\in\mathcal{P}_{2n+1}\) is the Hermite interpolant matching \(f\) and \(f'\) at each \(x_{j}\), the remainder is

\[
f(x)-H(x)
=
\frac{f^{(2n+2)}(\theta_{x})}{(2n+2)!}\,\omega(x)^{2}.
\]

Gauss integrates \(H\) exactly (\(\deg H\le 2n+1\)), and \(H(x_{j})=f(x_{j})\), so \(I_{n+1}(f)=I(H)\) and \(I(f)-I_{n+1}(f)=I(f-H)\). Integrating the remainder against the Chebyshev weight, and using \(I(T_{n+1}^{2})=\pi/2\), produces the printed bound

\[
\bigl|I(f)-I_{n+1}(f)\bigr)
\le
\frac{\pi\|f^{(2n+2)}\|_{\infty}}{(2n+2)!\,2^{2n+1}}.
\]

(The same \(\omega\) appears in (c) unsquared, which is why (c) has \((n+1)!\,2^{n}\) and a first-derivative gap of one, while (d) squares the nodal polynomial and doubles the order.)

\begin{examnote}[What ``derive \(I_{n+1}\)'' is]
\label{examnote:p2-quad}
The script for (d) has three lines: the nodes \eqref{eq:cheb-nodes}; the formula \eqref{eq:gauss-cheb}; and the error bound via \(\omega^{2}\) and \(I(T_{n+1}^{2})=\pi/2\). Naming ``Gauss--Chebyshev'' without writing the sum is not the formula. A common slip is to put \(1/\sqrt{1-x_{j}^{2}}\) into the sum: that factor is already inside each \(w_{j}\). Another is to use Chebyshev extrema \(\cos(j\pi/n)\) (Clenshaw--Curtis nodes) instead of the zeros \eqref{eq:cheb-nodes}. Part~(c) is the interpolant at those zeros, not the quadrature; the two share nodes and nothing else.
\end{examnote}

\subsubsection{Parts (a)--(c)}

Omitted. The closed form \(T_{n}(\cos\theta)=\cos(n\theta)\) from (a) is used below; the zeros of \(T_{n+1}\) are the nodes of (c), reused as Gauss nodes.

\subsubsection{Solution of Problem 2(d)}

\paragraph{Answer.}
Write \(N\coloneqq n+1\). The sum below is interpolatory quadrature at the zeros of \(T_{N}\); that is the definition of the \(N\)-node Gauss rule for \(I\) (\cref{rem:why-gauss}).

\begin{equation}
\label{eq:p2d-formula}
I_{N}(f)
=
\frac{\pi}{N}
\sum_{j=0}^{N-1}
f(x_{j}),
\qquad
x_{j}
=
\cos\theta_{j},
\qquad
\theta_{j}
=
\frac{(2j+1)\pi}{2N}.
\end{equation}

The nodes are the zeros of \(T_{N}\). The rule is exact on \(\mathcal{P}_{2N-1}=\mathcal{P}_{2n+1}\). For \(f\in C^{2n+2}[-1,1]\),

\[
\bigl|I(f)-I_{N}(f)\bigr|
\le
\frac{\pi\|f^{(2n+2)}\|_{\infty}}{(2n+2)!\,2^{2n+1}}.
\]

(The exam's \(I_{n+1}\) is this \(I_{N}\).)

\begin{remark}[Why this is Gauss, not a different interpolant]
\label{rem:why-gauss}
There is no separate object called ``Gauss interpolation''. Part~(c) is ordinary Lagrange interpolation at the zeros of \(T_{N}\). Part~(d) is a \emph{quadrature}: replace \(I(f)\) by a weighted sum of samples. Every such sum with \(N\) nodes can be made interpolatory in the sense of \cref{def:interpolatory}: take any distinct nodes, form the Lagrange interpolant \(L\in\mathcal{P}_{n}\), and set \(I_{N}(f)\coloneqq I(L)\). That rule is always exact on \(\mathcal{P}_{n}\), and that is all one gets for a generic choice of nodes.

Gauss is the same interpolatory construction with a constrained choice of nodes. Write \(\langle f,g\rangle_{w}\coloneqq I(fg)\). If the nodes are the zeros of a degree-\(N\) polynomial orthogonal to \(\mathcal{P}_{n}\) for this inner product, exactness doubles to \(\mathcal{P}_{2n+1}\) (\cref{def:gauss,lem:gauss-prec-cheb}): divide \(p=q\,T_{N}+s\) with \(\deg q,\deg s\le n\), the product \(qT_{N}\) integrates to zero against the weight, and the samples of \(p\) agree with those of \(s\). For the Chebyshev weight the orthogonal polynomials are \(\{T_{m}\}\), so the Gauss nodes are exactly the zeros of \(T_{N}\), i.e.\ the points in \eqref{eq:p2d-formula}. The equal weights \(\pi/N\) are then forced by uniqueness: there is only one interpolatory \(N\)-node rule exact on \(\mathcal{P}_{2n+1}\), and the proof below checks that this particular sum is that rule.

The later Hermite interpolant \(H\in\mathcal{P}_{2n+1}\) is used only for the error, not to define the formula. The formula itself samples \(f\), not \(f'\).
\end{remark}

\begin{lemma}[Hermite remainder at double nodes]
\label{lem:p2d-hermite}
Let \(x_{0},\ldots,x_{n}\) be distinct, write \(\omega(x)\coloneqq\prod_{j=0}^{n}(x-x_{j})\), and let \(H\in\mathcal{P}_{2n+1}\) satisfy \(H(x_{j})=f(x_{j})\) and \(H'(x_{j})=f'(x_{j})\) for each \(j\). If \(f\in C^{2n+2}[-1,1]\), then for every \(x\) there is \(\theta_{x}\) between \(x\) and the nodes such that

\[
f(x)-H(x)
=
\frac{f^{(2n+2)}(\theta_{x})}{(2n+2)!}\,\omega(x)^{2}.
\]
\end{lemma}

\begin{proof}
If \(x\) is a node both sides vanish. If not, choose \(K\) so that

\[
F(y)
\coloneqq
f(y)-H(y)-K\,\omega(y)^{2}
\]

vanishes at \(y=x\). Then \(F\) has a double zero at each \(x_{j}\) (because \(H\) matches \(f\) to first order and \(\omega^{2}\) has a double zero there) and a further zero at \(x\), hence at least \(2n+3\) zeros counting multiplicity. Rolle's theorem applied \(2n+2\) times produces \(\theta_{x}\) with \(F^{(2n+2)}(\theta_{x})=0\). Since \(\deg H\le 2n+1\) and \(\omega^{2}\) is monic of degree \(2n+2\), this is \(f^{(2n+2)}(\theta_{x})-K\cdot(2n+2)! =0\), so \(K=f^{(2n+2)}(\theta_{x})/(2n+2)!\). Evaluating \(F(x)=0\) is the claim.
\end{proof}

\begin{proof}[Solution of (d)]
\emph{The integral after the Chebyshev substitution.}
From (a), \(T_{m}(\cos\theta)=\cos(m\theta)\). The change of variables \(x=\cos\theta\), \(\theta\in[0,\pi]\), gives \(\mathrm{d}x/\sqrt{1-x^{2}}=-\mathrm{d}\theta\), hence

\[
I(f)
=
\int_{0}^{\pi} f(\cos\theta)\,\mathrm{d}\theta.
\]

In particular \(I(1)=\pi\) and, for every integer \(m\ge 1\),

\[
I(T_{m})
=
\int_{0}^{\pi}\cos(m\theta)\,\mathrm{d}\theta
=
0,
\qquad
I(T_{m}^{2})
=
\int_{0}^{\pi}\cos^{2}(m\theta)\,\mathrm{d}\theta
=
\int_{0}^{\pi}\frac{1+\cos(2m\theta)}{2}\,\mathrm{d}\theta
=
\frac{\pi}{2}.
\]

The zeros of \(T_{N}\) on \((-1,1)\) are exactly the points \(x_{j}\) in \eqref{eq:p2d-formula}: \(\cos(N\theta_{j})=\cos((2j+1)\pi/2)=0\).

\emph{The formula, by exhibiting a rule exact on \(\mathcal{P}_{2n+1}\).}
Define \(I_{N}\) by \eqref{eq:p2d-formula}. The family \(\{T_{0},\ldots,T_{2n+1}\}\) is a basis of \(\mathcal{P}_{2n+1}\), so it is enough to check \(I_{N}(T_{k})=I(T_{k})\) for \(k=0,\ldots,2n+1\). For \(k=0\),

\[
I_{N}(T_{0})
=
\frac{\pi}{N}\sum_{j=0}^{N-1}1
=
\pi
=
I(1).
\]

For \(1\le k\le 2n+1\), the target is \(I(T_{k})=0\), and

\[
I_{N}(T_{k})
=
\frac{\pi}{N}
\sum_{j=0}^{N-1}
\cos(k\theta_{j})
=
\frac{\pi}{N}\,
\operatorname{Re}
\sum_{j=0}^{N-1}
e^{ik\theta_{j}}.
\]

Write \(\zeta\coloneqq e^{ik\pi/N}\). Then

\[
\sum_{j=0}^{N-1}e^{ik\theta_{j}}
=
e^{ik\pi/(2N)}
\sum_{j=0}^{N-1}\zeta^{j}.
\]

One has \(\zeta=1\) if and only if \(k\) is a multiple of \(2N\). No such \(k\) occurs in \(\{1,\ldots,2N-1\}=\{1,\ldots,2n+1\}\), so

\[
\sum_{j=0}^{N-1}\zeta^{j}
=
\frac{1-\zeta^{N}}{1-\zeta}
=
\frac{1-(-1)^{k}}{1-\zeta}.
\]

If \(k\) is even the numerator vanishes and the sum is \(0\). If \(k\) is odd, write \(\varphi\coloneqq k\pi/N\). Then \(1-e^{i\varphi}=-2i\,e^{i\varphi/2}\sin(\varphi/2)\), and

\[
e^{i\varphi/2}\cdot\frac{2}{1-e^{i\varphi}}
=
\frac{i}{\sin(\varphi/2)}
\]

is purely imaginary. In both cases the real part vanishes, so \(I_{N}(T_{k})=0=I(T_{k})\). Thus \(I_{N}\) is exact on \(\mathcal{P}_{2n+1}\). An \((n+1)\)-node interpolatory rule exact on \(\mathcal{P}_{2n+1}\) is the Gauss rule for this weight (\cref{def:gauss,lem:gauss-prec-cheb}); the nodes are necessarily the zeros of \(T_{N}\).

\emph{The error bound.}
Let \(\omega(x)\coloneqq\prod_{j=0}^{n}(x-x_{j})\), the monic polynomial with the Gauss nodes as roots. The recurrence \(T_{m+1}=2xT_{m}-T_{m-1}\) with \(T_{1}(x)=x\) gives that the leading coefficient of \(T_{N}\) is \(2^{n}\), hence \(\omega=2^{-n}T_{N}\). Let \(H\in\mathcal{P}_{2n+1}\) be the Hermite interpolant of \(f\) at the double nodes \(x_{j}\). \Cref{lem:p2d-hermite} yields

\[
f(x)-H(x)
=
\frac{f^{(2n+2)}(\theta_{x})}{(2n+2)!}\,\omega(x)^{2}.
\]

Exactness on \(\mathcal{P}_{2n+1}\) gives \(I(H)=I_{N}(H)\). But \(H(x_{j})=f(x_{j})\), so \(I_{N}(H)=I_{N}(f)\). Therefore \(I(f)-I_{N}(f)=I(f-H)\). The weight \((1-x^{2})^{-1/2}\) is positive and \(\omega^{2}\ge 0\), so the pointwise remainder yields

\[
\bigl|I(f)-I_{N}(f)\bigr|
\le
\frac{\|f^{(2n+2)}\|_{\infty}}{(2n+2)!}\,I(\omega^{2}).
\]

The remaining factor is

\[
I(\omega^{2})
=
2^{-2n}I(T_{N}^{2})
=
2^{-2n}\cdot\frac{\pi}{2}
=
\frac{\pi}{2^{2n+1}},
\]

which is the printed constant.
\end{proof}

\begin{examnote}[What of this belongs on a script]
\label{examnote:p2d-script}
Write \eqref{eq:p2d-formula} first (nodes and the equal weight \(\pi/N\)). Exactness on \(\{T_{0},\ldots,T_{2n+1}\}\) is the derivation of the formula; the odd-\(k\) geometric sum must be checked to be purely imaginary, not merely quoted as ``the midpoint rule on \([0,\pi]\)''. The error is Hermite at double nodes, not the Lagrange remainder of (c): that is the only reason the factorial is \((2n+2)!\) and \(\omega\) is squared. Evaluating \(I(T_{N}^{2})=\pi/2\) is one line after \(x=\cos\theta\).
\end{examnote}


\subsection{Problem 3. [20 points]}
In many applications generalized eigenvalue problem is needed to be solved: finding \(x \in \mathbb{C}^{n}\) and \(\mu, \lambda \in \mathbb{C}\), s.t.\ \(x \neq 0\), \(|\mu|^{2} + |\lambda|^{2} = 1\), \(\mu A x = \lambda B x\) for given \(A, B \in \mathbb{C}^{n \times n}\). Here \(x\) and \((\lambda, \mu)\) are called the eigenvector and eigenvalue of the pair \((A, B)\) respectively.

If \(B\) is nonsingular, this problem is equivalent to the standard eigenvalue problem \(A B^{-1} y = \mu^{-1} \lambda y\).

For two unitary matrices \(Q, Z\), \(Q^{H} A Z = S\), \(Q^{H} B Z = T\) is called a UET performing on \((A, B)\), written as \(Q^{H}(A, B)Z = (S, T)\) for ease.
\begin{enumerate}[label=(\alph*)]
    \item Give an example that \((A, B)\) has more than \(n\) eigenvalues.
    \item Prove that there exists a UET such that \(Q^{H}(A, B)Z = (S, T)\), where \(S, T\) are upper triangular.
    \item Can you obtain the eigenvalues of \((A, B)\) from \((S, T)\)? How?
    \item Propose a method to construct a UET such that \(Q_{0}^{H}(A, B)Z_{0} = (A_{0}, B_{0})\), where \(A_{0}\) is upper Hessenberg, and \(B_{0}\) is upper triangular. For ease, you may assume \(n = 4\).
\end{enumerate}

\subsection{Problem 4. [15 points]}
Construct the Du Fort-Frankel scheme for the diffusion equation in 2D \(u_{t} = u_{xx} + u_{yy}\) and discuss its consistency and stability.

\setcounter{section}{4}
\setcounter{theorem}{0}
\setcounter{definition}{0}
\setcounter{remark}{0}
\setcounter{note}{0}

\subsubsection{Preliminaries: the three words the exam actually asks}

The PDE \(u_t=u_{xx}+u_{yy}\) is well-posed. The exam is not asking whether heat flows. It asks three things about a named difference scheme: write it, decide whether it is consistent with that PDE, and decide whether it is stable. The official CAM names for this circle are finite-difference methods, stability and convergence, and Lax equivalence \cite{gko95,leveque07,strikwerda04,laxrichtmyer56}. Finite elements \cite{brennerscott10} discretize the same Laplacian weakly and are a different language; they are not what ``construct the Du~Fort--Frankel scheme'' means.

\begin{center}
\begin{tikzpicture}[
  box/.style={draw, rounded corners=2pt, align=center, inner sep=5pt, font=\small},
  arr/.style={-{Latex}, thick}
]
  \node[box] (pde) at (0,2.35) {PDE \(u_t=\Delta u\)};
  \node[box] (sch) at (4.55,2.35) {difference scheme};
  \node[box] (con) at (0,0.55) {consistency\\(truncation \(\to 0\))};
  \node[box] (sta) at (4.55,0.55) {stability\\(no spurious growth)};
  \node[box] (cvg) at (2.275,-1.15) {convergence\\(\(U\to u\) as \(h,\tau\to 0\))};
  \draw[arr] (pde) -- (sch);
  \draw[arr] (sch) -- (con);
  \draw[arr] (sch) -- (sta);
  \draw[arr] (con) -- (cvg);
  \draw[arr] (sta) -- (cvg);
\end{tikzpicture}
\end{center}

Lax equivalence is the arrow into the bottom box: for a linear well-posed IVP, a consistent scheme converges if and only if it is stable. Qiuzhen almost never asks you to quote the theorem. It asks you to compute the two hypotheses. Du~Fort--Frankel is the standard pair in which those hypotheses split: the scheme is unconditionally von Neumann stable, and it is consistent with the heat equation only along certain mesh paths.

\paragraph{The objects on the grid.}

\begin{notation}[Two-dimensional heat grid]
\label{not:grid-p4}
Fix mesh sizes \(h_x,h_y>0\) and a time step \(\tau>0\). Write \(x_j=jh_x\), \(y_\ell=\ell h_y\), and \(t_n=n\tau\). The unknown \(U_{j,\ell}^n\) is a grid function intended to approximate \(u(x_j,y_\ell,t_n)\). When a single spatial step is enough one may take \(h_x=h_y=h\). The mesh ratios that appear in every heat-scheme calculation are
\[
\mu_x
\coloneqq
\frac{\tau}{h_x^2},
\qquad
\mu_y
\coloneqq
\frac{\tau}{h_y^2},
\qquad
\sigma
\coloneqq
\frac{1}{h_x^2}+\frac{1}{h_y^2}.
\]
The first two are parabolic CFL numbers; \(\sigma\) is the coefficient of the centre point in the five-point Laplacian. A scheme is \emph{explicit} if \(U^{n+1}\) is given by a formula in already-known time levels; it is \emph{implicit} if computing \(U^{n+1}\) requires solving a linear system that couples neighbouring nodes at time \(t_{n+1}\). Du~Fort--Frankel is three-level and explicit.
\end{notation}

Three objects that are easy to conflate, and that the next definitions keep separate:
\begin{enumerate}[label=\textup{(\roman*)}]
    \item \(U_{j,\ell}^n\): the grid function the scheme actually computes.
    \item \(u_{j,\ell}^n=u(x_j,y_\ell,t_n)\): samples of a smooth PDE solution. In general they do not satisfy the scheme.
    \item \(e_{j,\ell}^n=U_{j,\ell}^n-u_{j,\ell}^n\): the global error. Consistency is not a statement about \(e\).
\end{enumerate}

\begin{definition}[Centered second differences; five-point Laplacian]
\label{def:five-point}
On a uniform mesh in one variable,
\[
\delta_{xx}v_j
\coloneqq
\frac{v_{j+1}-2v_j+v_{j-1}}{h_x^2}.
\]
If \(v\) is smooth, Taylor expansion at \(x_j\) gives \(\delta_{xx}v=v_{xx}+\frac{1}{12}h_x^2 v_{xxxx}+O(h_x^4)\). Likewise \(\delta_{yy}\) in the \(y\)-variable. The standard five-point Laplacian on the plane is
\[
\Delta_h V_{j,\ell}
\coloneqq
\delta_{xx}V_{j,\ell}+\delta_{yy}V_{j,\ell}
=
\frac{V_{j+1,\ell}+V_{j-1,\ell}}{h_x^2}
+
\frac{V_{j,\ell+1}+V_{j,\ell-1}}{h_y^2}
-
2\sigma\, V_{j,\ell}.
\]
This is the spatial operator that FTCS, Richardson, and Du~Fort--Frankel all start from. The adjective ``centered'' means this stencil, not a one-sided second difference and not \((D^0)^2\).
\end{definition}

\paragraph{The named heat schemes, in the order that produces Du~Fort--Frankel.}

The exam assumes the 1D parent is known. The 2D scheme is the same replacement on \(\Delta_h\).

\begin{definition}[FTCS]
\label{def:ftcs}
The forward-time centered-space scheme for \(u_t=\Delta u\) is one-step and explicit:
\[
\frac{U_{j,\ell}^{n+1}-U_{j,\ell}^n}{\tau}
=
\Delta_h U_{j,\ell}^n.
\]
It is consistent of order \(O(\tau+h_x^2+h_y^2)\) with no mesh restriction, and von Neumann stable if and only if \(\mu_x+\mu_y\le 1/2\). That parabolic CFL is the reason one looks for an unconditionally stable alternative.
\end{definition}

\begin{definition}[Richardson / leapfrog heat]
\label{def:richardson}
The Richardson scheme keeps the centered Laplacian and centres the time difference as well:
\[
\frac{U_{j,\ell}^{n+1}-U_{j,\ell}^{n-1}}{2\tau}
=
\Delta_h U_{j,\ell}^n.
\]
The scheme is three-level (two-step): \(U^{n+1}\) depends on both \(U^n\) and \(U^{n-1}\). It is consistent of order \(O(\tau^2+h_x^2+h_y^2)\), and it is unconditionally von Neumann \emph{unstable}. The unstable mode is the one that sees the centre coefficient \(-2\sigma U^n\) as a source rather than as dissipation.
\end{definition}

\begin{definition}[Du~Fort--Frankel replacement]
\label{def:dff}
The Du~Fort--Frankel scheme is Richardson with the centre value \(U_{j,\ell}^n\) in \(\Delta_h\) replaced by the time average \(\frac12\bigl(U_{j,\ell}^{n+1}+U_{j,\ell}^{n-1}\bigr)\). Equivalently, the diagonal term \(2\sigma U_{j,\ell}^n\) is replaced by \(\sigma\bigl(U_{j,\ell}^{n+1}+U_{j,\ell}^{n-1}\bigr)\). In 1D this is the classical formula
\[
\frac{U_j^{n+1}-U_j^{n-1}}{2\tau}
=
\frac{U_{j+1}^n-U_j^{n+1}-U_j^{n-1}+U_{j-1}^n}{h^2}.
\]
In 2D, applying the same replacement in each coordinate, the scheme is
\[
\frac{U_{j,\ell}^{n+1}-U_{j,\ell}^{n-1}}{2\tau}
=
\frac{U_{j+1,\ell}^n-U_{j,\ell}^{n+1}-U_{j,\ell}^{n-1}+U_{j-1,\ell}^n}{h_x^2}
+
\frac{U_{j,\ell+1}^n-U_{j,\ell}^{n+1}-U_{j,\ell}^{n-1}+U_{j,\ell-1}^n}{h_y^2}.
\]
Collecting the unknown \(U_{j,\ell}^{n+1}\) on the left gives the explicit update
\[
\bigl(1+2\tau\sigma\bigr)\,U_{j,\ell}^{n+1}
=
\frac{2\tau}{h_x^2}\bigl(U_{j+1,\ell}^n+U_{j-1,\ell}^n\bigr)
+
\frac{2\tau}{h_y^2}\bigl(U_{j,\ell+1}^n+U_{j,\ell-1}^n\bigr)
+
\bigl(1-2\tau\sigma\bigr)\,U_{j,\ell}^{n-1}.
\]
The right-hand side uses only time levels \(n\) and \(n-1\). No linear system is solved: the scheme is explicit, even though it is three-level. Starting the march requires two initial time levels (e.g.\ one FTCS step from \(t=0\)).
\end{definition}

The algebraic identity that makes the truncation readable is the reason for the replacement, not a second scheme:
\[
\begin{aligned}
&U_{j+1,\ell}^n-U_{j,\ell}^{n+1}-U_{j,\ell}^{n-1}+U_{j-1,\ell}^n
\\
&\qquad=
\bigl(U_{j+1,\ell}^n-2U_{j,\ell}^n+U_{j-1,\ell}^n\bigr)
-
\bigl(U_{j,\ell}^{n+1}-2U_{j,\ell}^n+U_{j,\ell}^{n-1}\bigr),
\end{aligned}
\]
and likewise in \(y\). Du~Fort--Frankel is therefore Richardson minus a second time difference scaled by \(1/h_x^2\) and \(1/h_y^2\). That extra term is \(O\bigl((\tau/h_x)^2+(\tau/h_y)^2\bigr)\) on smooth solutions, which is the whole consistency story.

\paragraph{Consistency is a residual, not an error.}

\begin{definition}[Local truncation error]
\label{def:truncation}
Write a scheme in PDE form as an operator \(\mathcal L_{h,\tau}\) on grid functions, set equal to zero. For Du~Fort--Frankel that operator is the difference between the two sides of \cref{def:dff}. The computed unknown satisfies \(\mathcal L_{h,\tau}U=0\). Let \(u\) be a smooth solution of \(u_t=\Delta u\), and write \(u_{j,\ell}^n=u(x_j,y_\ell,t_n)\) for its samples. The \emph{local truncation error} (or residual) is
\[
\mathcal T_{j,\ell}^n
\coloneqq
\bigl(\mathcal L_{h,\tau}u\bigr)_{j,\ell}^n.
\]
This is an explicit number once \(u\) is given: every term is a value of a known function. The order is the largest \((p,q,r)\) for which \(\mathcal T=O(\tau^p+h_x^q+h_y^r)\) along the family of meshes under consideration. Consistency is the statement that \(\mathcal T\to 0\); it is not a statement about \(e=U-u\), and it is not a statement about an amplification factor \(|g|\).
\end{definition}

\begin{definition}[Consistency]
\label{def:consistency}
The scheme \(\mathcal L_{h,\tau}\) is \emph{consistent} with the PDE if, at every fixed \((x,y,t)\) in the space--time domain,
\[
\mathcal T_{j,\ell}^n
\to
0
\qquad
\text{as }
(h_x,h_y,\tau)\to(0,0,0),
\]
for every sufficiently smooth solution \(u\). Equivalently, in a discrete norm, \(\|\mathcal T^n\|_h\to 0\) uniformly for \(t_n\le T\).

The limit is taken along the family of meshes actually used. If a restriction such as \(\tau/h_x\to 0\) is part of that family, consistency is only claimed under that restriction. If \(h_x,h_y,\tau\to 0\) independently with \(\tau/h_x\) bounded away from zero, the residual of Du~Fort--Frankel need not vanish: the scheme is then consistent with a different PDE (a hyperbolic regularization of heat), not with \(u_t=\Delta u\). That mesh dependence is the reason the exam asks consistency and stability as two separate sentences.
\end{definition}

\paragraph{Stability is a bound on the computed grid function.}

\begin{definition}[Lax--Richtmyer stability]
\label{def:lax-richtmyer}
Write a linear scheme as a marching operator on grid functions (for a two-step scheme, on pairs \((U^n,U^{n-1})\)), with a discrete norm \(\|\cdot\|_h\) (usually \(\ell_2\)). The scheme is \emph{stable} on a time interval \([0,T]\) if there exist \(C,\alpha\), independent of the mesh for all sufficiently small steps (possibly subject to a restriction such as \(\mu_x+\mu_y\le 1/2\)), such that
\[
\|U^n\|_h
\le
C\,e^{\alpha t_n}\,\|U^0\|_h,
\qquad
0\le t_n\le T.
\]
What is forbidden is growth like \(|g|^n\) with \(|g|>1\) independent of \(\tau\): that explodes as the mesh is refined with \(t_n\) fixed. Growth of the form \(e^{\alpha T}\) is allowed, because the PDE itself may grow. For the heat equation on the line the PDE does not grow, and the practical test is \(\|U^n\|_h\le C\|U^0\|_h\) with \(C\) independent of \(h,\tau\).
\end{definition}

On the whole line or a torus, constant-coefficient schemes are diagonalized by spatial exponentials. Checking that no frequency grows is then necessary and sufficient for an \(\ell_2\) bound. That computation is von Neumann analysis; it is not a second definition of stability, and it is not a Fourier transform in time.

\begin{definition}[von Neumann / Fourier criterion]
\label{def:von-neumann}
A constant-coefficient linear scheme on \(\mathbb Z^2\) admits Fourier modes
\[
U_{j,\ell}^n
=
g^n\,e^{\mathrm i(j\xi+\ell\eta)},
\qquad
\xi,\eta\in\mathbb R.
\]
Here \((\xi,\eta)\) are dimensionless frequencies (one may write \(\xi=\kappa h_x\) with \(\kappa\) the dual variable to \(x\)). Substituting and cancelling the common factor \(g^{n-1}e^{\mathrm i(j\xi+\ell\eta)}\) produces a polynomial in the unknown amplification factor \(g\). For a one-step scheme the polynomial is linear and there is a unique \(g(\xi,\eta)\). For a two-step scheme (Du~Fort--Frankel, Richardson, and the wave scheme of Problem~5) the polynomial is quadratic,
\[
A(\xi,\eta)\,g^2+B(\xi,\eta)\,g+C(\xi,\eta)=0,
\]
with two roots \(g_\pm\). Both must be controlled: one is the physical mode, the other is a parasitic computational mode created by the extra time level.

The scheme is von Neumann stable if there is \(K\), independent of \(\xi,\eta\) and of the mesh, such that
\[
\bigl|g_\pm(\xi,\eta)\bigr|
\le
1+K\tau
\qquad
\text{for all }\xi,\eta.
\]
For heat (dissipative, no exponential PDE growth) this collapses to \(|g_\pm|\le 1\) for all frequencies. The GKO form \(1+K\tau\) is the general statement on the syllabus \cite{gko95}; a factor \(1+K\tau\) per step is \(e^{KT}\) at time \(T\), which is PDE growth, not mesh blow-up.
\end{definition}

\begin{definition}[Conditional / unconditional von Neumann stability]
\label{def:cond-uncond}
After Fourier substitution one has amplification factors \(g(\xi,\eta;h,\tau)\).
\begin{itemize}
    \item The scheme is von Neumann stable for a given mesh if \(|g|\le 1\) (or \(|g|\le 1+K\tau\)) for every frequency.
    \item It is \emph{conditionally stable} if that inequality holds only under an extra restriction, typically a bound on \(\mu=\tau/h^2\). FTCS is the model: stable iff \(\mu_x+\mu_y\le 1/2\).
    \item It is \emph{unconditionally stable} if the inequality holds for every \(h_x,h_y,\tau>0\). Backward Euler on heat is the model; Du~Fort--Frankel is the explicit three-level analogue.
    \item It is \emph{unconditionally unstable} if, for every mesh however fine, some frequency has \(|g|>1\). Richardson heat is the model. ``Unconditional'' here modifies the instability: there is no mesh on which the scheme becomes stable.
\end{itemize}
Consistency (\cref{def:consistency}) is a different question. A scheme may be consistent and unconditionally unstable (Richardson), or stable and inconsistent along some mesh paths (Du~Fort--Frankel). Neither adjective is implied by ``explicit'' or ``three-level''.
\end{definition}

\begin{proposition}[Jury's criterion for a real quadratic]
\label{prop:jury}
Let \(p(g)=Ag^2+Bg+C\) with \(A>0\) and \(A,B,C\in\mathbb R\). Both roots satisfy \(|g|\le 1\) if and only if
\[
p(1)\ge 0,
\qquad
p(-1)\ge 0,
\qquad
|C|\le A.
\]
Roots on the unit circle are permitted for von Neumann stability. This is the test that turns the Du~Fort--Frankel quadratic into three scalar inequalities in \((\tau,h_x,h_y,\xi,\eta)\).
\end{proposition}

\paragraph{Why the exam asks both.}

\begin{definition}[Convergence]
\label{def:convergence}
The scheme converges to a PDE solution \(u\) if, in the same discrete norm used for stability,
\[
\|U^n-u(\,\cdot\,,t_n)\|_h
\to
0
\qquad
\text{as }
(h_x,h_y,\tau)\to(0,0,0),
\]
uniformly for \(t_n\le T\), for every smooth initial datum (and, for two-step schemes, a consistent choice of the second starting level). Convergence is a statement about \(e=U-u\). Consistency is a statement about \(\mathcal T\). They are not synonyms.
\end{definition}

\begin{theorem}[Lax equivalence]
\label{thm:lax-eq}
Consider a linear well-posed evolutionary IVP and a consistent finite-difference approximation, in a discrete norm in which the PDE itself is well-posed. The scheme is convergent if and only if it is stable \cite{laxrichtmyer56,leveque07}.
\end{theorem}

The hypotheses are not decoration. Linearity produces a marching operator independent of the solution; well-posedness of the PDE makes the exact solution a legitimate comparison object; the same discrete norm is used for stability and for convergence. Von Neumann is the \(\ell_2\) machine. A discrete maximum principle (Problem~5) is the \(\ell_\infty\) machine.

In exam writing one checks consistency by Taylor expansion of \(\mathcal T\) and stability by von Neumann or Jury; Lax then gives convergence for free, \emph{along those mesh families for which both hypotheses hold}. If one refines Du~Fort--Frankel with \(\tau/h_x\) bounded away from zero, consistency with the heat equation fails, so Lax does not yield convergence to heat solutions, even though \(|g|\le 1\) for every \(\tau\). Unconditional stability is not a licence to take \(\tau\) comparable to \(h\).

\begin{remark}[What of this is the exam]
\label{rem:p4-script}
``Construct'' means write \cref{def:dff} and the explicit update. ``Consistency'' means expand \(\mathcal T\) and record the mesh restriction \(\tau/h_x\to 0\), \(\tau/h_y\to 0\). ``Stability'' means insert \(g^n e^{\mathrm i(j\xi+\ell\eta)}\), obtain the quadratic, and apply \cref{prop:jury} (or an equivalent \(|g_\pm|\le 1\) argument) for every \(\tau>0\). Lax equivalence is the warrant that those two computations decide convergence; it is not itself an answer.
\end{remark}

\subsubsection{Solution of Problem 4}

This is the two-dimensional Du~Fort--Frankel scheme for heat. It is not Problem~5 (the wave scheme with weights \(\frac14,\frac12,\frac14\)). The Fourier substitution below is the calculation on the handwritten page: a 2D mode \(g^n e^{\mathrm i(j\xi+\ell\eta)}\), and the centre value of \(\Delta_h\) replaced by \(\frac12(U^{n+1}+U^{n-1})\). The quadratic that appears is \emph{not} palindromic, so \(|g|=1\) is the wrong target. The target is \(|g_\pm|\le 1\) for every frequency and every mesh.

\paragraph{Construct.}

Start from Richardson (leapfrog time, centered Laplacian) for \(u_t=u_{xx}\),
\[
\frac{U_j^{n+1}-U_j^{n-1}}{2\tau}
=
\frac{U_{j+1}^n-2U_j^n+U_{j-1}^n}{h^2},
\]
and replace the centre value \(2U_j^n\) by \(U_j^{n+1}+U_j^{n-1}\). That is Du~Fort--Frankel in one dimension:
\[
\frac{U_j^{n+1}-U_j^{n-1}}{2\tau}
=
\frac{U_{j+1}^n-U_j^{n+1}-U_j^{n-1}+U_{j-1}^n}{h^2}.
\]
Do the same replacement in each coordinate of \(\Delta_h\). The two-dimensional scheme is
\[
\frac{U_{j,\ell}^{n+1}-U_{j,\ell}^{n-1}}{2\tau}
=
\frac{U_{j+1,\ell}^n-U_{j,\ell}^{n+1}-U_{j,\ell}^{n-1}+U_{j-1,\ell}^n}{h_x^2}
+
\frac{U_{j,\ell+1}^n-U_{j,\ell}^{n+1}-U_{j,\ell}^{n-1}+U_{j,\ell-1}^n}{h_y^2}.
\]
Write \(\sigma=1/h_x^2+1/h_y^2\). The unknown \(U_{j,\ell}^{n+1}\) appears linearly on both sides. Collecting it gives the explicit update
\[
\bigl(1+2\tau\sigma\bigr)\,U_{j,\ell}^{n+1}
=
\frac{2\tau}{h_x^2}\bigl(U_{j+1,\ell}^n+U_{j-1,\ell}^n\bigr)
+
\frac{2\tau}{h_y^2}\bigl(U_{j,\ell+1}^n+U_{j,\ell-1}^n\bigr)
+
\bigl(1-2\tau\sigma\bigr)\,U_{j,\ell}^{n-1}.
\]
The right-hand side uses only time levels \(n\) and \(n-1\). No spatial linear system is solved. The scheme is three-level and explicit. Two initial time levels are required to start.

\paragraph{Consistency.}

Feed a smooth solution \(u\) of \(u_t=u_{xx}+u_{yy}\) into the scheme. The local truncation error is the residual
\begin{align*}
\mathcal T_{j,\ell}^n
&=
\frac{u_{j,\ell}^{n+1}-u_{j,\ell}^{n-1}}{2\tau}
-
\frac{u_{j+1,\ell}^n-u_{j,\ell}^{n+1}-u_{j,\ell}^{n-1}+u_{j-1,\ell}^n}{h_x^2}
\\
&\qquad
-
\frac{u_{j,\ell+1}^n-u_{j,\ell}^{n+1}-u_{j,\ell}^{n-1}+u_{j,\ell-1}^n}{h_y^2}.
\end{align*}
Expand at \((x_j,y_\ell,t_n)\). The centered time difference is
\[
\frac{u(t_n+\tau)-u(t_n-\tau)}{2\tau}
=
u_t+\frac{\tau^2}{6}u_{ttt}+O(\tau^4)
=
u_t+O(\tau^2).
\]
(The even powers cancel; the first surviving correction is odd in \(\tau\), hence \(O(\tau^2)\) after dividing by \(2\tau\).)

For the \(x\)-term, expand the four samples separately. Odd powers of \(h_x\) from \(u_{j\pm 1}\) cancel, and odd powers of \(\tau\) from \(u^{n\pm 1}\) cancel, leaving
\begin{align*}
u_{j+1,\ell}^n+u_{j-1,\ell}^n-u_{j,\ell}^{n+1}-u_{j,\ell}^{n-1}
&=
\bigl(2u+h_x^2 u_{xx}+O(h_x^4)\bigr)
-
\bigl(2u+\tau^2 u_{tt}+O(\tau^4)\bigr)
\\
&=
h_x^2 u_{xx}-\tau^2 u_{tt}+O(h_x^4+\tau^4).
\end{align*}
Dividing by \(h_x^2\),
\[
\frac{u_{j+1,\ell}^n-u_{j,\ell}^{n+1}-u_{j,\ell}^{n-1}+u_{j-1,\ell}^n}{h_x^2}
=
u_{xx}-\Bigl(\frac{\tau}{h_x}\Bigr)^2 u_{tt}
+
O\Bigl(h_x^2+\frac{\tau^4}{h_x^2}\Bigr).
\]
The \(y\)-term is the same with \((h_y,u_{yy})\). Subtracting from the time difference,
\[
\mathcal T
=
\bigl(u_t-u_{xx}-u_{yy}\bigr)
+
\Biggl[\Bigl(\frac{\tau}{h_x}\Bigr)^2+\Bigl(\frac{\tau}{h_y}\Bigr)^2\Biggr]u_{tt}
+
O\Bigl(\tau^2+h_x^2+h_y^2+\frac{\tau^4}{h_x^2}+\frac{\tau^4}{h_y^2}\Bigr).
\]
On solutions of the heat equation the first parenthesis vanishes, so
\[
\mathcal T
=
\Biggl[\Bigl(\frac{\tau}{h_x}\Bigr)^2+\Bigl(\frac{\tau}{h_y}\Bigr)^2\Biggr]u_{tt}
+
O\Bigl(\tau^2+h_x^2+h_y^2+\frac{\tau^4}{h_x^2}+\frac{\tau^4}{h_y^2}\Bigr).
\]
Consistency means \(\mathcal T\to 0\) as \((h_x,h_y,\tau)\to(0,0,0)\) along the meshes actually used (\cref{def:consistency}). The displayed leading term does not vanish merely because each of \(h_x,h_y,\tau\) tends to zero: the path \(\tau=h_x=h_y\to 0\) leaves a remainder \(2u_{tt}\). The scheme is consistent with \(u_t=\Delta u\) if and only if, in addition,
\[
\frac{\tau}{h_x}\to 0,
\qquad
\frac{\tau}{h_y}\to 0.
\]
Under that restriction the remainder \(O(\tau^4/h_x^2)=\tau^2\cdot(\tau/h_x)^2\) also tends to zero. Parabolic refinement \(\tau=O(h_x^2)\) (and likewise in \(y\)) implies the restriction, hence consistency. If instead \(\tau/h_x\) stays bounded away from zero, the scheme is consistent with the hyperbolic regularization
\[
u_t
+
\Biggl[\Bigl(\frac{\tau}{h_x}\Bigr)^2+\Bigl(\frac{\tau}{h_y}\Bigr)^2\Biggr]u_{tt}
=
u_{xx}+u_{yy},
\]
not with the heat equation.

\paragraph{Stability.}

On \(\mathbb Z^2\) (or a torus) substitute the Fourier mode
\[
U_{j,\ell}^n
=
g^n e^{\mathrm i(j\xi+\ell\eta)},
\qquad
\xi,\eta\in\mathbb R.
\]
Every shift in \(j\) multiplies by \(e^{\mathrm i\xi}\), every shift in \(\ell\) multiplies by \(e^{\mathrm i\eta}\), and a time step multiplies by \(g\). Cancel the never-vanishing factor \(g^n e^{\mathrm i(j\xi+\ell\eta)}\). The left-hand side of the scheme becomes
\[
\frac{g-g^{-1}}{2\tau}.
\]
The \(x\)-contribution on the right is
\[
\frac{e^{\mathrm i\xi}-g-g^{-1}+e^{-\mathrm i\xi}}{h_x^2}
=
\frac{2\cos\xi}{h_x^2}-\frac{g+g^{-1}}{h_x^2},
\]
and likewise in \(y\). Adding,
\[
\frac{g-g^{-1}}{2\tau}
=
C-\sigma\bigl(g+g^{-1}\bigr),
\]
where
\[
C
\coloneqq
\frac{2\cos\xi}{h_x^2}+\frac{2\cos\eta}{h_y^2}.
\]
Because \(|\cos\xi|\le 1\) and \(|\cos\eta|\le 1\),
\[
|C|
\le
\frac{2}{h_x^2}+\frac{2}{h_y^2}
=
2\sigma.
\]
This bound replaces any further trigonometry.

Multiply the cancelled equation by \(g\) (equivalently: restore \(U^{n+1}=g\cdot U^n\) and \(U^{n-1}=g^{-1}U^n\)):
\[
\frac{g^2-1}{2\tau}
=
Cg-\sigma(g^2+1).
\]
Bring every term to the left:
\[
\Bigl(\frac{1}{2\tau}+\sigma\Bigr)g^2
-
C g
+
\Bigl(\sigma-\frac{1}{2\tau}\Bigr)
=
0.
\]
The coefficient of \(g^2\) is \(1/(2\tau)+\sigma\), not \(1/(2\tau)-\sigma\). The constant term is \(\sigma-1/(2\tau)\), not \(-(\sigma+1/(2\tau))\). Those two sign errors produce a different polynomial, and a modulus computed from that polynomial is not the amplification factor of the scheme.

Clear the remaining \(2\tau\) to obtain the equivalent quadratic
\[
(1+2\tau\sigma)\,g^2
-
2\tau C\,g
+
(2\tau\sigma-1)
=
0.
\]
Write \(A=1+2\tau\sigma>0\), \(B=-2\tau C\), \(C_{\mathrm{poly}}=2\tau\sigma-1\). The product of the roots is \(C_{\mathrm{poly}}/A=(2\tau\sigma-1)/(2\tau\sigma+1)\), which is not \(1\) in general. Heat dissipates: one must prove \(|g|\le 1\), not \(|g|=1\). Do not insert the quadratic formula and expand \(|C\pm\sqrt{\Delta}|/|2A|\).

Apply Jury's criterion (\cref{prop:jury}) to \(p(g)=Ag^2+Bg+C_{\mathrm{poly}}\). Both roots satisfy \(|g|\le 1\) if and only if \(p(1)\ge 0\), \(p(-1)\ge 0\), and \(|C_{\mathrm{poly}}|\le A\).

\[
p(1)
=
A+B+C_{\mathrm{poly}}
=
4\tau\sigma-2\tau C
=
2\tau(2\sigma-C).
\]
The bound \(C\le 2\sigma\) gives \(p(1)\ge 0\), with equality at \((\xi,\eta)=(0,0)\).

\[
p(-1)
=
A-B+C_{\mathrm{poly}}
=
4\tau\sigma+2\tau C
=
2\tau(2\sigma+C).
\]
The bound \(C\ge -2\sigma\) gives \(p(-1)\ge 0\).

Finally \(|2\tau\sigma-1|\le 1+2\tau\sigma\). If \(2\tau\sigma\ge 1\) this is \(2\tau\sigma-1\le 1+2\tau\sigma\), i.e.\ \(-1\le 1\). If \(2\tau\sigma<1\) this is \(1-2\tau\sigma\le 1+2\tau\sigma\), i.e.\ \(0\le 4\tau\sigma\). Both hold for every \(\tau\sigma>0\).

The three Jury inequalities are therefore valid for every \(\xi,\eta\) and every \(\tau,h_x,h_y>0\). The scheme is unconditionally von Neumann stable.

When the discriminant is negative the roots are conjugate and
\[
|g|^2
=
\frac{|2\tau\sigma-1|}{1+2\tau\sigma}
<1
\qquad(\tau\sigma>0),
\]
which is dissipation, as expected for heat. When the discriminant is nonnegative the roots are real and Jury already keeps them in \([-1,1]\). Either way, the explicit square root is unnecessary.

\paragraph{Summary.}

The two-dimensional Du~Fort--Frankel scheme is explicit and unconditionally von Neumann stable. It is consistent with \(u_t=u_{xx}+u_{yy}\) if and only if \(\tau/h_x\to 0\) and \(\tau/h_y\to 0\). Unconditional stability does not cancel that mesh restriction. Along meshes with \(\tau/h\) bounded away from zero, Lax equivalence (\cref{thm:lax-eq}) does not yield convergence to heat solutions, because consistency fails.

\begin{examnote}[Script for Problem 4]
\label{examnote:p4}
Construct: write the 2D replacement and the explicit factor \(1+2\tau\sigma\). Consistency: Taylor to \(\mathcal T=[(\tau/h_x)^2+(\tau/h_y)^2]u_{tt}+O(\cdots)\), then \(\tau/h_x\to 0\), \(\tau/h_y\to 0\). Stability: the cancelled cosine identity, the quadratic \((1+2\tau\sigma)g^2-2\tau C g+(2\tau\sigma-1)=0\) with \(|C|\le 2\sigma\), then Jury. Do not take moduli of the quadratic formula. Do not quote \(|g|=1\); that is Problem~5.
\end{examnote}


\subsection{Problem 5. [15 points]}
For the wave equation \(u_{tt} = u_{xx}\), analyze the stability of the scheme

\[
\frac{u_{j}^{n+1} - 2u_{j}^{n} + u_{j}^{n-1}}{\tau^{2}}
=
\frac{u_{j+1}^{n+1} - 2u_{j}^{n+1} + u_{j-1}^{n+1}}{4h^{2}}
+
\frac{u_{j+1}^{n} - 2u_{j}^{n} + u_{j-1}^{n}}{2h^{2}}
+
\frac{u_{j+1}^{n-1} - 2u_{j}^{n-1} + u_{j-1}^{n-1}}{4h^{2}}.
\]

\setcounter{section}{5}
\setcounter{theorem}{0}
\setcounter{definition}{0}
\setcounter{remark}{0}
\setcounter{note}{0}

\subsubsection{What ``analyze the stability'' means}

The PDE \(u_{tt}=u_{xx}\) is the one-dimensional wave equation. On the line it is well-posed: energy \(\int(u_t^2+u_x^2)\) is conserved, and solutions oscillate rather than explode. The exam is not asking whether the PDE is stable. It asks whether \emph{this difference scheme} can manufacture exponentially growing grid functions that the PDE does not have.

The machine is von Neumann analysis, already defined in Problem~4 (\cref{def:von-neumann,def:lax-richtmyer,def:cond-uncond}). The calculation below does not assume that those definitions have been memorized: every algebraic step is written out. The target is a single inequality. Insert a Fourier mode \(U_j^n=g^n e^{\mathrm i j\xi}\). After one time step the mode is multiplied by a complex number \(g=g(\xi;\tau,h)\). After \(n\) steps it is multiplied by \(g^n\). If some frequency has \(|g|>1\), then \(|g|^{t_n/\tau}\) explodes as the mesh is refined with \(t_n\) fixed, and the scheme is unstable. If \(|g|\le 1\) for every \(\xi\) and every \(\tau,h>0\), the scheme is unconditionally von Neumann stable.

The scheme is three-level, so \(g\) solves a quadratic and there are two roots \(g_\pm\). Both must satisfy \(|g|\le 1\). One is the physical oscillation; the other is a parasitic computational mode created by the extra time level.

Consistency is not asked. Do not Taylor-expand the residual.

\paragraph{Step 0. Rewrite the stencil.}

The three fractions on the right-hand side are the same centered second difference, evaluated at three consecutive times, with weights \(\frac14,\frac12,\frac14\). Write
\[
\delta_{xx}V_j
\coloneqq
\frac{V_{j+1}-2V_j+V_{j-1}}{h^2}.
\]
The scheme is
\[
\frac{U_j^{n+1}-2U_j^n+U_j^{n-1}}{\tau^2}
=
\frac14\,\delta_{xx}U_j^{n+1}
+
\frac12\,\delta_{xx}U_j^n
+
\frac14\,\delta_{xx}U_j^{n-1}.
\]
The unknown \(U^{n+1}\) appears on both sides, so a linear system in the spatial index \(j\) is solved at each step: the scheme is implicit in space, two-step in time.

\paragraph{Step 1. Substitute a Fourier mode.}

On the infinite line (or a periodic interval) the shift \(j\mapsto j\pm 1\) is diagonalized by complex exponentials. It is therefore enough to test the scheme on
\[
U_j^n
=
g^n e^{\mathrm i j\xi},
\qquad
\xi\in\mathbb R,\quad g\in\mathbb C.
\]
The number \(\xi\) is a dimensionless spatial frequency. The number \(g\) is the unknown amplification factor: it is the ratio of the mode at time \(t_{n+1}\) to the same mode at time \(t_n\). Different frequencies may have different \(g\); that is why one writes \(g(\xi)\). The calculation must hold for \emph{every} \(\xi\), not just for \(\xi=0\) or \(\xi=\pi\).

Write the eight terms that appear. The left-hand side is
\[
\frac{g^{n+1}e^{\mathrm i j\xi}-2g^n e^{\mathrm i j\xi}+g^{n-1}e^{\mathrm i j\xi}}{\tau^2}
=
g^{n-1}e^{\mathrm i j\xi}\cdot\frac{g^2-2g+1}{\tau^2}
=
g^{n-1}e^{\mathrm i j\xi}\cdot\frac{(g-1)^2}{\tau^2}.
\]
Each spatial second difference acts on a pure exponential, so it factors. This is the step that, if skipped, produces a page of complex moduli.

\paragraph{Step 2. The spatial symbol (do not write \((e^{\mathrm i j\xi}-1)^2\)).}

Compute the centered second difference of \(e^{\mathrm i j\xi}\) by shifting the index, not by a formal square:
\begin{align*}
e^{\mathrm i(j+1)\xi}-2e^{\mathrm i j\xi}+e^{\mathrm i(j-1)\xi}
&=
e^{\mathrm i j\xi}\bigl(e^{\mathrm i\xi}-2+e^{-\mathrm i\xi}\bigr)
\\
&=
e^{\mathrm i j\xi}\bigl(2\cos\xi-2\bigr)
\\
&=
-4\sin^2(\xi/2)\,e^{\mathrm i j\xi}.
\end{align*}
The last line is \(1-\cos\xi=2\sin^2(\xi/2)\). In particular the result is a \emph{real} multiple of \(e^{\mathrm i j\xi}\), and the multiple does not depend on the grid point \(j\).

The expression \((e^{\mathrm i j\xi}-1)^2=e^{2\mathrm i j\xi}-2e^{\mathrm i j\xi}+1\) is a different polynomial. It still contains \(j\), so it cannot be the symbol of a translation-invariant stencil. Replacing \(\delta_{xx}\) by that square is the error that forces one to compute \(|B\pm\sqrt{\Delta}|/|2A|\) with complex \(A,B\).

Dividing by \(h^2\),
\[
\delta_{xx}\bigl(e^{\mathrm i j\xi}\bigr)
=
\lambda(\xi)\,e^{\mathrm i j\xi},
\qquad
\lambda(\xi)
\coloneqq
-\frac{4\sin^2(\xi/2)}{h^2}
\le 0.
\]
The same factor \(\lambda(\xi)\) multiplies the mode at every time level. The right-hand side of the scheme is therefore
\[
\lambda(\xi)\,g^{n-1}e^{\mathrm i j\xi}
\left(
\frac14 g^2+\frac12 g+\frac14
\right)
=
\lambda(\xi)\,g^{n-1}e^{\mathrm i j\xi}\cdot\frac{(g+1)^2}{4}.
\]

\paragraph{Step 3. Cancel the common factor.}

The mode \(g^{n-1}e^{\mathrm i j\xi}\) is never zero. Cancel it from both sides:
\[
\frac{(g-1)^2}{\tau^2}
=
\frac{\lambda(\xi)}{4}(g+1)^2.
\]
The coefficients are now ordinary complex numbers built from \(\tau,h,\xi\) only. There is no \(j\) left, and no uncancelled exponential.

Introduce the nonnegative mesh quantity
\[
q
\coloneqq
\frac{\tau^2}{h^2}\sin^2\frac{\xi}{2}
\ge 0.
\]
Then \(\lambda(\xi)\tau^2/4=-q\), and the cancelled equation is
\[
(g-1)^2
=
-q\,(g+1)^2.
\]
If \(g=0\), the left side is \(1\) and the right side is \(-q\le 0\), so \(1=-q\) is impossible. If \(g=-1\), the left side is \(4\) and the right side is \(0\), so \(4=0\) is impossible. Thus \(g\notin\{0,-1\}\), and one may divide.

\paragraph{Step 4. Reduce to \(g+1/g\).}

Divide \((g-1)^2=-q(g+1)^2\) by \(g\):
\[
\frac{g-2+g^{-1}}{\tau^2}
=
\frac{\lambda(\xi)}{4}\bigl(g+2+g^{-1}\bigr),
\]
which is the same as
\[
g+\frac{1}{g}-2
=
-q\left(g+2+\frac{1}{g}\right).
\]
Collect the terms that contain \(g+1/g\):
\[
\left(1+q\right)\left(g+\frac{1}{g}\right)
=
2-2q,
\]
hence
\[
g+\frac{1}{g}
=
2\frac{1-q}{1+q}.
\]
Call the right-hand side \(z\). For every \(q\ge 0\) one has \(-1\le(1-q)/(1+q)\le 1\), so
\[
z
=
2\frac{1-q}{1+q}
\in[-2,2].
\]
This interval membership is the whole stability proof. It uses only \(q\ge 0\), which holds for every \(\tau,h>0\) and every \(\xi\). No restriction of CFL type appears.

\paragraph{Step 5. A real number \(z\in[-2,2]\) forces \(|g|=1\).}

The identity \(g+1/g=z\) rearranges to the quadratic
\[
g^2-z g+1=0.
\]
The product of the two roots is \(1\). The discriminant is \(z^2-4\le 0\) because \(|z|\le 2\). Hence the roots are complex conjugates (or a double real root \(\pm 1\) when \(z=\pm 2\)). A pair of conjugate numbers whose product is \(1\) lies on the unit circle:
\[
|g|^2
=
g\,\overline{g}
=
1.
\]
Equivalently, the quadratic formula gives
\[
g
=
\frac{z\pm\mathrm i\sqrt{4-z^2}}{2},
\]
and
\[
|g|^2
=
\frac{z^2+(4-z^2)}{4}
=
1.
\]
Both roots satisfy \(|g|=1\). This is von Neumann stability, and it is stronger than the inequality \(|g|\le 1\): every Fourier mode of the scheme oscillates at constant amplitude, matching the conservative character of \(u_{tt}=u_{xx}\).

The same conclusion in trigonometric language: since \(z\in[-2,2]\), there is a real angle \(\varphi\) with \(\cos\varphi=z/2=(1-q)/(1+q)\), and \(g=e^{\pm\mathrm i\varphi}\).

\paragraph{Step 6. Unconditional stability.}

The argument never restricted \(\tau/h\). For the explicit centered leapfrog scheme
\[
\frac{U_j^{n+1}-2U_j^n+U_j^{n-1}}{\tau^2}
=
\delta_{xx}U_j^n,
\]
the analogous computation produces \(g+1/g=2-4(\tau/h)^2\sin^2(\xi/2)\), which exits \([-2,2]\) as soon as \(\tau/h>1\). The \(\frac14,\frac12,\frac14\) average is what keeps \(z\) inside \([-2,2]\) for every mesh: the implicit Laplacian at time \(n+1\) supplies enough dissipation in the amplification polynomial to kill the CFL barrier, without introducing amplitude decay (\(|g|\) stays exactly \(1\), not \(<1\)).

\begin{theorem}[Von Neumann stability of the exam scheme]
\label{thm:p5-stable}
The scheme is unconditionally von Neumann stable: for every \(\tau,h>0\) and every frequency \(\xi\), both amplification factors satisfy \(|g_\pm(\xi)|=1\).
\end{theorem}

\paragraph{The quadratic you were trying to write, with real coefficients.}

After Step~3, move every term to the left:
\[
\left(\frac{1}{\tau^2}-\frac{\lambda}{4}\right)g^2
-
\left(\frac{2}{\tau^2}+\frac{\lambda}{2}\right)g
+
\left(\frac{1}{\tau^2}-\frac{\lambda}{4}\right)
=
0.
\]
In the \(q\)-notation this is
\[
(1+q)\,g^2-2(1-q)\,g+(1+q)=0.
\]
Write \(A=1+q>0\), \(B=-2(1-q)\), \(C=A\). These are real. The palindromic shape \(A=C\) is exactly \(g\cdot(1/g)=1\), which is the product-of-roots observation on the handwritten page. The discriminant is
\[
\Delta
=
4(1-q)^2-4(1+q)^2
=
-16q
\le 0.
\]
For real palindromic quadratics, \(\Delta\le 0\) plus \(C/A=1\) is \(|g|=1\), which is Step~5 again. One does not insert these \(A,B\) into
\[
g=\frac{-B\pm\sqrt{\Delta}}{2A}
\]
and then take moduli of complex numerators and denominators. That expansion is valid but pointless: \(\Delta\le 0\) has already placed the square root on the imaginary axis, and \(A=C\) has already placed the roots on the unit circle.

Jury's criterion (\cref{prop:jury}) is a third packaging of the same three inequalities \(p(1)\ge 0\), \(p(-1)\ge 0\), \(|C|\le A\), all true for every \(q\ge 0\). On a script, Step~4 plus Step~5 is shorter.

\begin{remark}[What the two roots are]
\label{rem:p5-two-roots}
A two-step scheme needs two starting time levels, and its symbol is quadratic, so two modes live at each \(\xi\). The root with \(g\to 1\) as \(\xi\to 0\) is the physical wave (it approximates \(e^{\pm\mathrm i(\tau/h)\xi}\) to second order). The second root is parasitic: it is an artefact of writing a second-order equation as a three-level recurrence. Because \(|g|=1\) for that root as well, the parasite does not grow. It can still pollute the solution if the second initial level is chosen inconsistently; that is a question about starting the march, not about von Neumann stability of the recurrence.
\end{remark}

\begin{remark}[What this does not prove]
\label{rem:p5-not-proved}
Stability is not consistency, and it is not convergence. The scheme is in fact consistent of order \(O(\tau^2+h^2)\) with \(u_{tt}=u_{xx}\) (Taylor-expand at \((x_j,t_n)\); the \(\frac14,\frac12,\frac14\) weights are the Simpson average of \(\delta_{xx}\) in time, and no \(\tau^2/h^2\) remainder of Du~Fort--Frankel type appears). Combined with \cref{thm:p5-stable} and Lax equivalence (\cref{thm:lax-eq}), one would get convergence. The exam did not ask for that sentence. Phase error (the discrete \(\varphi\) versus the PDE frequency) is also not a stability question: \(|g|=1\) says nothing about whether the numerical wave speed is accurate.
\end{remark}

\begin{examnote}[What belongs on a script]
\label{examnote:p5}
Four lines. (i)~Rewrite as \(\frac14,\frac12,\frac14\) on \(\delta_{xx}\). (ii)~Substitute \(g^n e^{\mathrm i j\xi}\) and use \(\delta_{xx}e^{\mathrm i j\xi}=-4\sin^2(\xi/2)\,e^{\mathrm i j\xi}/h^2\). (iii)~Cancel, set \(q=(\tau/h)^2\sin^2(\xi/2)\), obtain \(g+1/g=2(1-q)/(1+q)\in[-2,2]\). (iv)~Conclude \(|g_\pm|=1\) for every \(\tau,h,\xi\). Do not expand \(|B\pm\sqrt{\Delta}|/|2A|\). Do not replace the second difference by \((e^{\mathrm i j\xi}-1)^2\).
\end{examnote}


\noindent\textbf{You can choose either 6 or 7 to answer. The points will be decided as \(\max(6, 7)\).}

\subsection{Problem 6. [20 points]}
Consider the so-called Rayleigh oscillator

\[
y'' + y + \epsilon\Bigl[\tfrac13 (y')^{3} - y'\Bigr] = 0,
\]

with initial condition

\[
y(0) = 0,\qquad y'(0) = 2a,
\]

where \(y = y(t)\), \(\epsilon > 0\), \(a > 0\).
\begin{enumerate}[label=(\alph*)]
    \item (14 points) For a small \(\epsilon\), construct an approximation of the solution to the above problem which is valid for large \(t\).
    \item (4 points) What is the accuracy of this approxmation? State a conclusion and briefly explain it.
    \item (2 points) Plot the approximated orbits in the phase plane, i.e., \(y - y'\) plane for several different \(a\). And explain what you observe.
\end{enumerate}

\setcounter{section}{6}
\setcounter{theorem}{0}
\setcounter{definition}{0}
\setcounter{remark}{0}
\setcounter{note}{0}

\subsubsection{Preliminaries, then the script}

Rewrite the Rayleigh equation as the weakly forced oscillator
\begin{equation}
\label{eq:p6-std}
\ddot y+y
=
\epsilon\Bigl(\dot y-\tfrac13(\dot y)^3\Bigr),
\qquad
y(0)=0,\quad \dot y(0)=2a,
\end{equation}
with \(\epsilon>0\) small and \(a>0\) fixed. If \(\epsilon=0\), the solution is the circle
\begin{equation}
\label{eq:p6-linear}
y(t)=2a\sin t,
\qquad
\dot y(t)=2a\cos t
\end{equation}
of radius \(2a\) in the \((y,\dot y)\) plane. Part~(a) asks for an approximation that remains ordered when \(t\) is large, meaning \(t\) of size \(1/\epsilon\), not merely \(t=O(1)\). The force is size \(\epsilon\) at each instant; integrated over a time \(1/\epsilon\) it produces an \(O(1)\) change in amplitude. A regular series in the single variable \(t\) cannot see that change. Two-time multiple scales can.

Two facts are used throughout. The first says the amplitude must move. The second says why a term \(t\sin t\) cannot be tolerated in a first correction.

\begin{lemma}[Energy]
\label{lem:p6-energy}
Along solutions of \eqref{eq:p6-std},
\begin{equation}
\label{eq:p6-energy}
\frac{\mathrm{d}E}{\mathrm{d}t}
=
\epsilon\,\dot y^2\Bigl(1-\tfrac13\dot y^2\Bigr),
\qquad
E
\coloneqq
\tfrac12\dot y^2+\tfrac12 y^2.
\end{equation}
Thus \(E\) increases wherever \(\lvert\dot y\rvert<\sqrt{3}\) and decreases wherever \(\lvert\dot y\rvert>\sqrt{3}\).
\end{lemma}

\begin{proof}
Multiply \eqref{eq:p6-std} by \(\dot y\). The left-hand side is \(\frac{\mathrm{d}}{\mathrm{d}t}(\frac12\dot y^2+\frac12 y^2)\). The right-hand side is \(\epsilon(\dot y^2-\frac13\dot y^4)\).
\end{proof}

On a nearly circular orbit of radius \(R\), one has \(\lvert\dot y\rvert\) of size \(R\). Small circles grow and large circles shrink, so the amplitude cannot stay at the initial value \(2a\). The time for an \(O(1)\) change is \(T=\epsilon t\), because \(\mathrm{d}E/\mathrm{d}t=O(\epsilon)\).

\begin{lemma}[Resonant forcing]
\label{lem:p6-resonant}
The equation \(\ddot u+u=\cos 3t\) has the bounded particular solution \(-\frac18\cos 3t\). The equation \(\ddot u+u=\cos t\) has the particular solution \(\frac12 t\sin t\), and no bounded particular solution of the form \(C\cos t+D\sin t\). The same holds with \(\sin\) in place of \(\cos\), giving \(-\frac12 t\cos t\).
\end{lemma}

\begin{proof}
Direct differentiation: \(\frac{\mathrm{d}^2}{\mathrm{d}t^2}(\cos 3t)+\cos 3t=-8\cos 3t\), and \(\frac{\mathrm{d}^2}{\mathrm{d}t^2}(t\sin t)+t\sin t=2\cos t\). A guess \(C\cos t\) lies in the kernel of \(\frac{\mathrm{d}^2}{\mathrm{d}t^2}+1\), so it cannot match \(\cos t\).
\end{proof}

The factor \(t\) in \cref{lem:p6-resonant} is the secular term. If it appears in \(y_1\), then \(\epsilon y_1\) is \(O(1)\) already at \(t\sim 1/\epsilon\), and the expansion has stopped being ordered.

\subsubsection{(a) An approximation valid for large \(t\)}

\paragraph{The regular expansion, which is not the answer.}

Substitute \(y=y_0(t)+\epsilon y_1(t)+\cdots\) into \eqref{eq:p6-std}, with \(y_0(0)=0\), \(\dot y_0(0)=2a\), and homogeneous data on \(y_1\). Order \(\epsilon^0\) is \(\ddot y_0+y_0=0\), hence
\[
y_0(t)=2a\sin t,
\qquad
\dot y_0(t)=2a\cos t.
\]
Order \(\epsilon^1\) is
\[
\ddot y_1+y_1
=
\dot y_0-\tfrac13(\dot y_0)^3
=
2a\cos t-\tfrac83 a^3\cos^3 t.
\]
The identity \(\cos^3 t=(3\cos t+\cos 3t)/4\) rewrites the right-hand side as
\begin{equation}
\label{eq:p6-y1-rhs}
2a(1-a^2)\cos t-\tfrac23 a^3\cos 3t.
\end{equation}
By \cref{lem:p6-resonant},
\begin{equation}
\label{eq:p6-y1}
y_1(t)
=
a(1-a^2)\,t\sin t
+
\text{(bounded harmonics)}.
\end{equation}
On any interval of \(\epsilon\)-independent length the product \(\epsilon t\sin t\) is \(O(\epsilon)\), and \(y_0+\epsilon y_1\) is consistent. At \(t=c/\epsilon\) the same product is \(O(1)\). The regular expansion is therefore not an approximation valid for large \(t\). (If \(a=1\) the resonant coefficient vanishes and \(y_0=2\sin t\) is already the attractor; that is a special initial value, not the argument for general \(a>0\).)

\begin{remark}[Why the regular expansion is not the answer]
\label{rem:p6-regular-fails}
Yes: it is not valid for the large times the exam asks about. The phrase ``valid for large \(t\)'' in part~(a) means a time of size \(1/\epsilon\), not an \(\epsilon\)-independent interval. Write \(T=\epsilon t\) for the slow clock that appears in the next paragraph. Then
\[
t=O(1)
\qquad\Longleftrightarrow\qquad
T=O(\epsilon),
\qquad
t=\frac{c}{\epsilon}
\qquad\Longleftrightarrow\qquad
T=c=O(1).
\]
The regular series is ordered on the first of those regimes and unordered on the second. It is the second regime that the exam wants. (``Large \(T\)'' in the sense \(T\to\infty\) is later still, \(t\sim 1/\epsilon^2\); that is the horizon of part~(b), not the reason \(y_0+\epsilon y_1\) is discarded.)

Two equivalent readings of the same failure.

\emph{Asymptotic.} The correction is \(\epsilon y_1=\epsilon a(1-a^2)\,t\sin t+\cdots\). A first-order expansion is legitimate only while \(\epsilon y_1=o(y_0)\). On \(t=O(1)\) one has \(\epsilon t=O(\epsilon)\), so the product is a genuine \(O(\epsilon)\) correction. At \(t=c/\epsilon\) one has \(\epsilon t=c\), so
\[
\epsilon y_1
=
a(1-a^2)\,c\sin t+O(\epsilon).
\]
The ``small'' term is now \(O(1)\), the same size as \(y_0=2a\sin t\). The series \(y_0+\epsilon y_1+\cdots\) has stopped being ordered: the first correction is no longer smaller than the leading term. That is what secular means. The bounded harmonics in \eqref{eq:p6-y1} are harmless; only the factor \(t\) is fatal.

\emph{Energetic.} \Cref{lem:p6-energy} says \(\mathrm{d}E/\mathrm{d}t=O(\epsilon)\), so the radius must change by \(O(1)\) once \(t\) has length \(1/\epsilon\). The profile \(y_0=2a\sin t\) freezes the radius at the initial value \(2a\) forever. A uniform-in-\(t\) approximation cannot do that: small circles must grow and large circles must shrink. The regular expansion never moves the amplitude, so it cannot remain close to the true orbit on that interval.

The case \(a=1\) is not a rescue. The resonant coefficient in \eqref{eq:p6-y1-rhs} vanishes, \(y_1\) stays bounded, and \(2\sin t\) happens to be the limit cycle. That is one initial datum. Part~(a) is for general \(a>0\), for which the secular term is present and the radius must run from \(2a\) to \(2\).
\end{remark}

\paragraph{Two-time expansion.}

The secular product \(\epsilon t\sin t\) is \(T\sin t\) with \(T=\epsilon t\). Absorb it into a slowly varying amplitude by allowing the leading profile to depend on both clocks. Write
\begin{equation}
\label{eq:p6-ms-ansatz}
y(t)=Y_0(t,T)+\epsilon Y_1(t,T)+\cdots,
\qquad
T=\epsilon t,
\end{equation}
and treat \(t\) and \(T\) as independent arguments until the end. The chain rule on a composite \(Y(t,\epsilon t)\) is
\begin{equation}
\label{eq:p6-chain}
\frac{\mathrm{d}}{\mathrm{d}t}=\partial_t+\epsilon\partial_T,
\qquad
\frac{\mathrm{d}^2}{\mathrm{d}t^2}=\partial_{tt}+2\epsilon\partial_{tT}+\epsilon^2\partial_{TT}.
\end{equation}
Substitute into \eqref{eq:p6-std}.

\emph{Order \(\epsilon^0\).} \(\partial_{tt}Y_0+Y_0=0\). The general solution may still depend on the unused variable \(T\):
\begin{equation}
\label{eq:p6-polar}
Y_0(t,T)=R(T)\cos\bigl(t+\varphi(T)\bigr).
\end{equation}
Write \(\theta\coloneqq t+\varphi(T)\). If \(R\) and \(\varphi\) were constant one would recover the regular expansion.

\emph{Order \(\epsilon^1\).} Differentiating \eqref{eq:p6-polar} at fixed \(T\) gives \(\partial_t Y_0=-R\sin\theta\) and
\[
\partial_{tT}Y_0=-R'\sin\theta-R\varphi'\cos\theta.
\]
The \(O(\epsilon)\) balance is
\begin{equation}
\label{eq:p6-Y1}
\partial_{tt}Y_1+Y_1
=
-2\partial_{tT}Y_0+\partial_t Y_0-\tfrac13(\partial_t Y_0)^3.
\end{equation}
The first summand comes from \eqref{eq:p6-chain}; the remaining two are the leading piece of \(\dot y-\frac13\dot y^3\). Using \((\partial_t Y_0)^3=-R^3\sin^3\theta\) and \(\sin^3\theta=(3\sin\theta-\sin 3\theta)/4\), the right-hand side of \eqref{eq:p6-Y1} is
\begin{equation}
\label{eq:p6-Y1-rhs}
\bigl(2R'-R+\tfrac14 R^3\bigr)\sin\theta
+
2R\varphi'\cos\theta
-
\tfrac{1}{12}R^3\sin 3\theta.
\end{equation}
The \(\sin 3\theta\) term is nonresonant and produces a bounded piece of \(Y_1\). The \(\sin\theta\) and \(\cos\theta\) terms are resonant. If either coefficient is nonzero, \cref{lem:p6-resonant} puts a factor \(t\) into \(Y_1\), and the same failure as \eqref{eq:p6-y1} reappears. Setting both coefficients to zero is therefore a design condition on \(R\) and \(\varphi\):
\begin{equation}
\label{eq:p6-amp}
R'
=
\frac{R}{2}\Bigl(1-\frac{R^2}{4}\Bigr)
=
\frac{R}{8}(4-R^2),
\qquad
\varphi'=0
\quad(R\neq 0).
\end{equation}
The first equation is the slow energy balance in polar coordinates. The second says there is no \(O(\epsilon)\) frequency correction.

\paragraph{Initial data and the explicit amplitude.}

The equilibria of \(R'=R(4-R^2)/8\) are \(R=0\) and \(R=\pm 2\). The derivative of the right-hand side is \(\frac12-\frac{3R^2}{8}\), equal to \(+\frac12\) at \(0\) and to \(-1\) at \(2\), so \(R=0\) is unstable and \(R=2\) is attracting on \(R>0\).

At \((t,T)=(0,0)\),
\[
Y_0(0,0)=R(0)\cos\varphi(0)=0,
\qquad
\partial_t Y_0(0,0)=-R(0)\sin\varphi(0)=2a.
\]
With \(a>0\) and \(R>0\) this forces \(\varphi(0)=-\pi/2\) and \(R(0)=2a\). Since \(\varphi'=0\),
\[
Y_0(t,T)=R(T)\sin t.
\]
Separate variables in \eqref{eq:p6-amp}:
\[
\frac{8\,\mathrm{d}R}{R(4-R^2)}=\mathrm{d}T,
\qquad
\frac{1}{R(4-R^2)}
=
\frac14\Bigl(\frac{1}{R}+\frac{R}{4-R^2}\Bigr).
\]
Integration gives \(\ln(R^2/\lvert 4-R^2\rvert)=T+C\), or \(R^2/(4-R^2)=Ke^{T}\). The value \(R(0)=2a\) fixes \(K=a^2/(1-a^2)\). Solving for \(R\),
\begin{equation}
\label{eq:p6-R}
R(T)
=
\frac{2a}{\sqrt{a^2+(1-a^2)e^{-T}}}.
\end{equation}
(The formula remains valid for \(a>1\): the denominator \(a^2+(1-a^2)e^{-T}=a^2(1-e^{-T})+e^{-T}\) stays positive.) As \(T\to\infty\), \(R(T)\to 2\) for every \(a>0\). The approximation requested in (a) is therefore
\begin{equation}
\label{eq:p6-approx}
y(t)
\sim
R(\epsilon t)\sin t
=
\frac{2a\sin t}{\sqrt{a^2+(1-a^2)e^{-\epsilon t}}}.
\end{equation}

\begin{warning}[What not to write]
\label{warn:p6-lindstedt-vs-ms}
Quoting a limit-cycle theorem in place of \eqref{eq:p6-R} identifies the attractor \(R=2\) and misses the transient from radius \(2a\). Lindstedt--Poincar\'e stretches the clock to correct a period at fixed amplitude; energy is not conserved here, so the amplitude must run. The 2025 Spring twin \(\ddot y+y+\epsilon(\dot y)^3=0\) is the same machine with a decaying \(R\).
\end{warning}

\begin{examnote}[Part~(a) on a script]
\label{examnote:p6-script}
Write \eqref{eq:p6-energy} in one line. Write the regular expansion through \eqref{eq:p6-y1} and the size count at \(t\sim 1/\epsilon\). Then write \eqref{eq:p6-ms-ansatz}--\eqref{eq:p6-polar}, the chain rule, the two coefficients in \eqref{eq:p6-Y1-rhs} set to zero, \(R(0)=2a\), \(\varphi(0)=-\pi/2\), and \eqref{eq:p6-approx}.
\end{examnote}

\subsubsection{(b) Accuracy}

After one multiple-scale step,
\begin{equation}
\label{eq:p6-acc}
y(t)=R(\epsilon t)\sin t+O(\epsilon)
\qquad\text{uniformly for}\qquad
0\le t\le\frac{C}{\epsilon},
\end{equation}
for any fixed \(C>0\). Equivalently, the error is \(O(\epsilon)\) on the interval where the slow time \(T\) remains \(O(1)\).

The count is the same one that killed the regular expansion, applied to a residual that is no longer resonant. The conditions \eqref{eq:p6-amp} remove every \(t\sin\theta\) and \(t\cos\theta\) from \(Y_1\). The first omitted terms in \eqref{eq:p6-chain} and in the expansion of the force are \(O(\epsilon^2)\) while \(T=O(1)\). Integrating an \(O(\epsilon^2)\) residual over a time of length \(O(1/\epsilon)\) produces an \(O(\epsilon)\) error. The regular expansion \eqref{eq:p6-y1} still has a resonant \(O(\epsilon)\) residual, so the same integral is \(O(1)\) on that interval.

The claim is not \(O(\epsilon)\) for all \(t\ge 0\). At \(t\sim 1/\epsilon^2\) one has \(T\sim 1/\epsilon\), a new secular obstruction appears, and another slow time would be required. That is the accuracy statement 2025 Spring asked for on the damped twin.

\subsubsection{(c) Orbits in the \((y,\dot y)\) plane}

Differentiating \eqref{eq:p6-approx} at leading order (the \(T\)-derivative of the envelope is \(O(\epsilon)\) and is dropped) gives \(\dot y\sim R(\epsilon t)\cos t\). Eliminating \(t\),
\[
y^2+\dot y^2\sim R(\epsilon t)^2.
\]
The approximated orbit is a circle whose radius changes slowly from \(R(0)=2a\) to \(2\). If \(0<a<1\) the orbit starts inside the circle of radius \(2\) and spirals out; if \(a>1\) it starts outside and spirals in; if \(a=1\) it sits on the attractor already. Every \(a>0\) forgets the initial radius on the time scale \(t\sim 1/\epsilon\).

\begin{center}
\begin{tikzpicture}[scale=0.72]
  \draw[-{Latex}] (-3.4,0) -- (3.6,0) node[right] {\(y\)};
  \draw[-{Latex}] (0,-3.1) -- (0,3.3) node[above] {\(y'\)};
  \draw[thick] (0,0) circle (2);
  \draw[dashed] (0,0) circle (1.1);
  \draw[dashed] (0,0) circle (2.85);
  \node[font=\footnotesize] at (2.15,2.15) {\(R=2\)};
  \node[font=\footnotesize] at (0.55,1.35) {\(a<1\)};
  \node[font=\footnotesize] at (2.05,2.85) {\(a>1\)};
  \draw[-{Latex}] (1.1,0) arc (0:40:1.1);
  \draw[-{Latex}] (2.85,0) arc (0:35:2.85);
\end{tikzpicture}
\end{center}

\begin{summary}[Problem~6]
\label{sum:p6}
The regular expansion \(y=2a\sin t+\epsilon a(1-a^2)\,t\sin t+\cdots\) fails at \(t\sim 1/\epsilon\). Multiple scales with \(T=\epsilon t\) produces \(y\sim R(\epsilon t)\sin t\) and the explicit radius \eqref{eq:p6-R}. The error is \(O(\epsilon)\) on \(t=O(1/\epsilon)\). Every \(a>0\) is drawn to the circle of radius \(2\).
\end{summary}


\subsection{Problem 7. [20 points]}
Consider the minimization

\[
\min_{x \in \mathbb{R}^{n}} \|Ax - b\|_{2} + \gamma \|x\|_{1},
\]

with \(A \in \mathbb{R}^{m \times n}\), \(b \in \mathbb{R}^{m}\), and \(\gamma > 0\).
\begin{enumerate}[label=(\alph*)]
    \item (8 points) Derive the Lagrange dual of the equivalent problem:
    \[
    \min_{x,y} \|y\|_{2} + \gamma \|x\|_{1}, \quad \text{s.t. } Ax - b = y,
    \]
    with \(x \in \mathbb{R}^{n}\) and \(y \in \mathbb{R}^{m}\).
    \item (8 points) Suppose \(Ax^{*} - b \neq 0\) where \(x^{*}\) is an optimal point. Define \(r = \dfrac{Ax^{*} - b}{\|Ax^{*} - b\|_{2}}\). Show that
    \[
    \|A^{\top} r\|_{\infty} \le \gamma, \qquad r^{\top} A x^{*} + \gamma \|x^{*}\|_{1} = 0.
    \]
    \item (4 points) Show that if \(\|a_{i}\|_{2} < \gamma\) where \(a_{i}\) is the \(i\)-th column of \(A\), then \(x^{*}_{i} = 0\).
\end{enumerate}

\setcounter{section}{7}
\setcounter{theorem}{0}
\setcounter{definition}{0}
\setcounter{remark}{0}
\setcounter{note}{0}

\subsubsection{Preliminaries, then the script}

The problem is unconstrained convex minimization of
\begin{equation}
\label{eq:p7-primal}
f(x)
\coloneqq
\|Ax-b\|_{2}+\gamma\|x\|_{1},
\qquad
x\in\mathbb{R}^{n}.
\end{equation}
The \(\ell_1\) term is not differentiable on the coordinate hyperplanes, and \(\|Ax-b\|_{2}\) is not differentiable on \(\{Ax=b\}\). Part~(a) therefore rewrites the residual as an equality constraint and passes to the Lagrange dual, which is a constrained problem in a multiplier \(\nu\in\mathbb{R}^{m}\) with only ball constraints. Parts~(b) and~(c) are first-order optimality for \(f\) itself. The dictionary is short; each item is used on the next page.

\begin{definition}[The three norms]
\label{def:p7-lp}
\[
\|x\|_{1}=\sum_{i=1}^{n}|x_{i}|,
\qquad
\|x\|_{2}=\Bigl(\sum_{i=1}^{n}x_{i}^{2}\Bigr)^{1/2},
\qquad
\|x\|_{\infty}=\max_{i}|x_{i}|.
\]
The only two pairings needed below are Cauchy--Schwarz and H\"older:
\begin{equation}
\label{eq:p7-holder}
\lvert\nu^{\top}y\rvert\le\|\nu\|_{2}\|y\|_{2},
\qquad
\lvert u^{\top}x\rvert
=
\Bigl\lvert\sum_{i}u_{i}x_{i}\Bigr\rvert
\le
\sum_{i}\lvert u_{i}\rvert\lvert x_{i}\rvert
\le
\|u\|_{\infty}\|x\|_{1}.
\end{equation}
If \(a_{i}\) is the \(i\)-th column of \(A\), then \((A^{\top}r)_{i}=a_{i}^{\top}r\).
\end{definition}

\begin{definition}[Lagrangian and Lagrange dual]
\label{def:p7-dual}
For \(\min_{z}F(z)\) subject to \(h(z)=0\), the Lagrangian is \(L(z,\nu)=F(z)+\nu^{\top}h(z)\). The dual function is \(g(\nu)=\inf_{z}L(z,\nu)\), and the Lagrange dual problem is \(\sup_{\nu}g(\nu)\). Weak duality is the comparison \(g(\nu)\le F(z)\) for every feasible \(z\) and every \(\nu\). Here \(z=(x,y)\), \(F=\|y\|_{2}+\gamma\|x\|_{1}\), and the exam's constraint \(Ax-b=y\) is written \(h(x,y)=Ax-b-y\).
\end{definition}

\begin{lemma}[When \(\inf_{z}(\alpha\|z\|-u^{\top}z)\) is finite]
\label{lem:p7-norm-inf}
Let \(\alpha>0\). Then
\[
\inf_{y}\bigl(\|y\|_{2}-\nu^{\top}y\bigr)
=
\begin{cases}
0, & \|\nu\|_{2}\le 1,\\
-\infty, & \|\nu\|_{2}>1,
\end{cases}
\qquad
\inf_{x}\bigl(\gamma\|x\|_{1}+u^{\top}x\bigr)
=
\begin{cases}
0, & \|u\|_{\infty}\le\gamma,\\
-\infty, & \|u\|_{\infty}>\gamma.
\end{cases}
\]
\end{lemma}

\begin{proof}
If \(\|\nu\|_{2}\le 1\), then \(\|y\|_{2}-\nu^{\top}y\ge\|y\|_{2}-\|\nu\|_{2}\|y\|_{2}\ge 0\), and the value \(0\) is attained at \(y=0\). If \(\|\nu\|_{2}>1\), the ray \(y=t\nu\) with \(t\to+\infty\) gives \(t\|\nu\|_{2}-t\|\nu\|_{2}^{2}=t\|\nu\|_{2}(1-\|\nu\|_{2})\to-\infty\).

If \(\|u\|_{\infty}\le\gamma\), then \(\gamma\|x\|_{1}+u^{\top}x\ge\gamma\|x\|_{1}-\|u\|_{\infty}\|x\|_{1}\ge 0\), attained at \(x=0\). If \(\|u\|_{\infty}>\gamma\), some coordinate satisfies \(\lvert u_{j}\rvert>\gamma\). The ray \(x=-t\,\operatorname{sign}(u_{j})\,e_{j}\) gives \(\gamma t-t\lvert u_{j}\rvert=t(\gamma-\lvert u_{j}\rvert)\to-\infty\).
\end{proof}

\subsubsection{(a) The Lagrange dual}

The original \(f(x)\) and the constrained problem in the prompt are the same object: given any \(x\), the only feasible \(y\) is \(y=Ax-b\), and the objective becomes \(f(x)\). The reason for introducing \(y\) is that the two norms then act on separate variables, so the dual function splits.

Form the Lagrangian of the written constraint \(Ax-b-y=0\):
\begin{equation}
\label{eq:p7-L}
L(x,y,\nu)
=
\|y\|_{2}+\gamma\|x\|_{1}+\nu^{\top}(Ax-b-y).
\end{equation}
The dual function is the unconstrained infimum
\[
g(\nu)
=
\inf_{x\in\mathbb{R}^{n},\,y\in\mathbb{R}^{m}}L(x,y,\nu).
\]
There is no \(x\)--\(y\) product in \(L\) except through the already-fixed vector \(\nu\), so the infimum factors:
\begin{equation}
\label{eq:p7-g-split}
g(\nu)
=
\inf_{y}\bigl(\|y\|_{2}-\nu^{\top}y\bigr)
+
\inf_{x}\bigl(\gamma\|x\|_{1}+(A^{\top}\nu)^{\top}x\bigr)
-
b^{\top}\nu.
\end{equation}
Apply \cref{lem:p7-norm-inf} to each infimum (the second with \(u=A^{\top}\nu\)). Both are \(0\) precisely when
\[
\|\nu\|_{2}\le 1
\qquad\text{and}\qquad
\|A^{\top}\nu\|_{\infty}\le\gamma;
\]
otherwise at least one infimum is \(-\infty\), hence \(g(\nu)=-\infty\). On that set one is left with the linear function \(-b^{\top}\nu\). The Lagrange dual is therefore the explicit program
\begin{equation}
\label{eq:p7-dual-prog}
\sup_{\nu\in\mathbb{R}^{m}}\bigl(-b^{\top}\nu\bigr)
\qquad\text{subject to}\qquad
\|\nu\|_{2}\le 1,
\quad
\|A^{\top}\nu\|_{\infty}\le\gamma.
\end{equation}
Equivalently, after the change of sign \(\lambda=-\nu\) (which is the Lagrangian of the flipped writing \(y-Ax+b=0\)),
\begin{equation}
\label{eq:p7-dual-flip}
\sup_{\lambda\in\mathbb{R}^{m}}\,b^{\top}\lambda
\qquad\text{subject to}\qquad
\|\lambda\|_{2}\le 1,
\quad
\|A^{\top}\lambda\|_{\infty}\le\gamma.
\end{equation}
Either display is a correct answer. Weak duality is \(g(\nu)\le f(x)\) for every \(x\) and every feasible \(\nu\) (\cref{rem:p7-weak}). Strong duality holds because the constraint is affine and a feasible pair always exists (pick any \(x\), set \(y=Ax-b\)): the optimal values of \eqref{eq:p7-primal} and \eqref{eq:p7-dual-prog} coincide (\cref{rem:p7-strong}).

\begin{remark}[Weak duality]
\label{rem:p7-weak}
Write
\[
p^{*}
\coloneqq
\inf_{x\in\mathbb{R}^{n}}f(x)
=
\inf_{\substack{x,y\\Ax-b=y}}
\bigl(\|y\|_{2}+\gamma\|x\|_{1}\bigr)
\]
for the primal value, and
\[
d^{*}
\coloneqq
\sup_{\nu\in\mathbb{R}^{m}}g(\nu)
\]
for the dual value. Weak duality is the comparison of these two numbers through a pointwise inequality that does not mention optima.

Take any \(x\in\mathbb{R}^{n}\) and set \(y=Ax-b\), so that \((x,y)\) is feasible for the written constraint. The Lagrangian \eqref{eq:p7-L} then collapses to the primal objective:
\[
L(x,y,\nu)
=
\|y\|_{2}+\gamma\|x\|_{1}+\nu^{\top}(Ax-b-y)
=
f(x).
\]
The dual function is an unconstrained infimum over a larger set, hence cannot exceed any particular value of \(L\):
\[
g(\nu)
=
\inf_{x',y'}L(x',y',\nu)
\le
L(x,y,\nu)
=
f(x).
\]
This holds for every \(x\) and every multiplier \(\nu\), with no convexity assumption and with no restriction that \(\nu\) be dual-feasible \cite[\S5.1.3]{boyd04}. Taking the infimum in \(x\) on the right and the supremum in \(\nu\) on the left yields
\[
d^{*}\le p^{*}.
\]
The difference \(p^{*}-d^{*}\ge 0\) is the duality gap.

If \(\nu\) fails the two ball constraints of \eqref{eq:p7-dual-prog}, then \(g(\nu)=-\infty\) by \cref{lem:p7-norm-inf}, and the inequality is vacuous. If \(\nu\) is dual-feasible, then \(g(\nu)=-b^{\top}\nu\), and weak duality becomes a concrete numerical bound: every such \(\nu\) lower-bounds every Lasso value,
\[
-b^{\top}\nu
\le
f(x)
\qquad
\text{for all }x\in\mathbb{R}^{n}.
\]
In particular \(-b^{\top}\nu\le p^{*}\). That is the whole point of writing a dual even if one never proves that the gap is zero: any feasible \(\nu\) certifies a lower bound, and a better \(\nu\) certifies a tighter one.

Weak duality does not say that \(g(\nu)=f(x)\) at some pair, and it does not say that a minimizer of \(f\) exists. It is only the one-sided comparison \(g\le f\).
\end{remark}

\begin{remark}[Strong duality]
\label{rem:p7-strong}
Strong duality is the statement that the gap vanishes:
\[
d^{*}=p^{*}.
\]
The best lower bound produced by the dual is then exactly the primal value, so the two programs in \eqref{eq:p7-primal} and \eqref{eq:p7-dual-prog} have the same optimal value. This is not implied by weak duality, and it is not implied by convexity of \(f\) alone.

A standard sufficient condition is Slater's constraint qualification for convex programs. In the form needed here \cite[\S5.2.3]{boyd04}: the objective \(F(x,y)=\|y\|_{2}+\gamma\|x\|_{1}\) is convex, the only constraint \(Ax-b-y=0\) is affine, and there exists a feasible point in the relative interior of \(\operatorname{dom} F\). Under those hypotheses one has \(d^{*}=p^{*}\).

Each hypothesis is immediate. The function \(F\) is a sum of norms, hence convex, and it is finite at every \((x,y)\in\mathbb{R}^{n}\times\mathbb{R}^{m}\), so \(\operatorname{ri}(\operatorname{dom} F)\) is the whole space. The constraint is affine by inspection. A feasible pair always exists: pick any \(x\) and set \(y=Ax-b\). Slater therefore applies, and the optimal values of \eqref{eq:p7-primal} and \eqref{eq:p7-dual-prog} coincide.

(The same conclusion is Fenchel--Rockafellar applied to \(\|A\,\cdot\,-b\|_{2}+\gamma\|\cdot\|_{1}\): the conjugate of \(\|\cdot\|_{2}\) is the indicator of the Euclidean unit ball, the conjugate of \(\gamma\|\cdot\|_{1}\) is the indicator of the \(\ell_{\infty}\) ball of radius \(\gamma\), and both domains are the whole space.)

Attainment is a separate claim. Strong duality equates two numbers; it does not by itself produce a minimizer \(x^{*}\) or a maximizer \(\nu^{*}\). Both happen to exist here. The primal attains because \(f\) is continuous and coercive: \(\gamma>0\) gives \(f(x)\ge\gamma\|x\|_{1}\to\infty\) as \(\|x\|\to\infty\). The dual attains because \(-b^{\top}\nu\) is continuous and the feasible set
\[
\bigl\{\nu:\|\nu\|_{2}\le 1\bigr\}
\cap
\bigl\{\nu:\|A^{\top}\nu\|_{\infty}\le\gamma\bigr\}
\]
is compact (closed subset of a Euclidean ball). Thus there exist \(x^{*}\) and \(\nu^{*}\) with
\[
f(x^{*})=g(\nu^{*})=p^{*}=d^{*}.
\]
A pair with zero gap is a saddle of \(L\): writing \(y^{*}=Ax^{*}-b\),
\[
L(x^{*},y^{*},\nu)
\le
L(x^{*},y^{*},\nu^{*})
\le
L(x,y,\nu^{*})
\]
for every \((x,y)\) and every \(\nu\). The left inequality is primal feasibility (the constraint term vanishes at \((x^{*},y^{*})\), independently of \(\nu\)); the right inequality says that \(\nu^{*}\) attains the infimum that defines \(g\). That saddle reading is the route of \cref{rem:p7-from-dual}. Part~(b) does not need it.

Two statements that are easy to swap. Weak duality is \(g(\nu)\le f(x)\) for every pair, including non-optimal ones. Strong duality is \(d^{*}=p^{*}\), i.e.\ equality of the two optimal values, not equality of \(g\) and \(f\) at a random feasible pair. Equality \(g(\nu)=f(x)\) holds if and only if \(x\) is primal-optimal and \(\nu\) is dual-optimal (and then the gap at that pair is automatically zero).
\end{remark}

\begin{examnote}[Part~(a) on a script]
\label{examnote:p7-dual}
Write \eqref{eq:p7-L} and \eqref{eq:p7-g-split}. Evaluate the two infima by Cauchy--Schwarz and H\"older as in \cref{lem:p7-norm-inf} (the two rays that produce \(-\infty\) must be exhibited). State \eqref{eq:p7-dual-prog} or \eqref{eq:p7-dual-flip}. One sentence on weak duality (\(g\le f\)) and one on Slater (affine constraint, pick \(y=Ax-b\)) is enough; \cref{rem:p7-weak,rem:p7-strong} are not the script. Naming ``the Lagrange dual'' without the constraints on \(\nu\) is not a derivation.
\end{examnote}

\subsubsection{(b) First-order conditions at \(x^{*}\)}

Let \(x^{*}\) minimize \(f\), and assume \(y^{*}\coloneqq Ax^{*}-b\neq 0\). The exam's vector
\[
r
=
\frac{y^{*}}{\|y^{*}\|_{2}}
\]
is the Euclidean unit vector in the direction of the residual. It is also the gradient of \(\|\cdot\|_{2}\) at \(y^{*}\), which is the only reason it appears.

Because \(x^{*}\) is an unconstrained minimizer, every directional derivative of \(f\) at \(x^{*}\) is nonnegative: for every \(v\in\mathbb{R}^{n}\),
\begin{equation}
\label{eq:p7-dir}
\liminf_{t\downarrow 0}\frac{f(x^{*}+tv)-f(x^{*})}{t}
\ge 0.
\end{equation}
Compute the two pieces of \(f\) separately.

\emph{The residual term.} Write \(y(t)=y^{*}+tAv\). Then
\[
\|y(t)\|_{2}^{2}
=
\|y^{*}\|_{2}^{2}+2t\,(y^{*})^{\top}Av+t^{2}\|Av\|_{2}^{2},
\]
and the first-order expansion of the square root about \(\|y^{*}\|_{2}>0\) is
\begin{equation}
\label{eq:p7-res-diff}
\|y^{*}+tAv\|_{2}
=
\|y^{*}\|_{2}+t\,r^{\top}Av+o(t).
\end{equation}
(This is the ordinary derivative of \(\|\cdot\|_{2}\) at a nonzero point: \(\nabla\|y\|_{2}=y/\|y\|_{2}\).)

\emph{The \(\ell_1\) term.} For all sufficiently small \(\lvert t\rvert\),
\begin{equation}
\label{eq:p7-l1-diff}
\|x^{*}+tv\|_{1}-\|x^{*}\|_{1}
=
t\sum_{x_{i}^{*}\neq 0}\operatorname{sign}(x_{i}^{*})\,v_{i}
+
\lvert t\rvert\sum_{x_{i}^{*}=0}\lvert v_{i}\rvert.
\end{equation}
(On a coordinate with \(x_{i}^{*}\neq 0\), \(\lvert x_{i}^{*}+tv_{i}\rvert=\lvert x_{i}^{*}\rvert+t\operatorname{sign}(x_{i}^{*})v_{i}\) for small \(t\). On a coordinate with \(x_{i}^{*}=0\), one has \(\lvert tv_{i}\rvert\).)

\emph{Coordinatewise tests.} Take \(v=e_{i}\) and \(v=-e_{i}\). Write \(u\coloneqq A^{\top}r\), so that \(r^{\top}Ae_{i}=u_{i}=a_{i}^{\top}r\).

If \(x_{i}^{*}\neq 0\), both directions are admissible and \eqref{eq:p7-l1-diff} is linear in \(t\). Condition \eqref{eq:p7-dir} in both directions forces the derivative to vanish:
\begin{equation}
\label{eq:p7-active}
u_{i}+\gamma\operatorname{sign}(x_{i}^{*})=0,
\qquad\text{i.e.}\qquad
\lvert u_{i}\rvert=\gamma.
\end{equation}

If \(x_{i}^{*}=0\), the \(\ell_1\) contribution is \(\gamma\lvert t\rvert\). The right derivative of \(f\) along \(+e_{i}\) is \(u_{i}+\gamma\), and the right derivative along \(-e_{i}\) is \(-u_{i}+\gamma\). Both are \(\ge 0\), hence
\begin{equation}
\label{eq:p7-inactive}
\lvert u_{i}\rvert\le\gamma.
\end{equation}

Combining \eqref{eq:p7-active} and \eqref{eq:p7-inactive} over every coordinate,
\begin{equation}
\label{eq:p7-inf}
\|A^{\top}r\|_{\infty}
=
\|u\|_{\infty}
\le\gamma,
\end{equation}
which is the first claim. For the second claim, pair \(u\) against \(x^{*}\). On every coordinate with \(x_{i}^{*}=0\) the product \(u_{i}x_{i}^{*}\) is zero. On every coordinate with \(x_{i}^{*}\neq 0\), equation \eqref{eq:p7-active} gives \(u_{i}x_{i}^{*}=-\gamma\lvert x_{i}^{*}\rvert\). Summing,
\begin{equation}
\label{eq:p7-comp}
r^{\top}Ax^{*}+\gamma\|x^{*}\|_{1}
=
u^{\top}x^{*}+\gamma\|x^{*}\|_{1}
=
0.
\end{equation}

\begin{remark}[The same argument in subdifferential language]
\label{rem:p7-subgrad}
Convexity of \(f\) says that \(x^{*}\) minimizes \(f\) if and only if \(0\in\partial f(x^{*})\). The chain rule and the sum rule give
\[
\partial f(x^{*})
=
A^{\top}\partial\|\cdot\|_{2}(y^{*})+\gamma\,\partial\|x^{*}\|_{1}.
\]
For \(y^{*}\neq 0\) one has \(\partial\|y^{*}\|_{2}=\{r\}\). The set \(\partial\|x\|_{1}\) is exactly the collection of vectors \(s\) with \(\|s\|_{\infty}\le 1\) and \(s_{i}=\operatorname{sign}(x_{i})\) whenever \(x_{i}\neq 0\); in particular \(s^{\top}x=\|x\|_{1}\). Thus \(A^{\top}r+\gamma s=0\) for some such \(s\), which is \eqref{eq:p7-inf}--\eqref{eq:p7-comp} again. This is 2024 Fall Problem~7(ii). The directional-derivative argument above does not require the vocabulary.
\end{remark}

\begin{remark}[From the dual of (a)]
\label{rem:p7-from-dual}
The same two displays are dual feasibility plus zero gap. Differentiability of \(\|y\|_{2}\) at \(y^{*}\neq 0\) forces the optimal multiplier of \eqref{eq:p7-L} to be \(\nu^{*}=r\). Dual feasibility of \(r\) is \(\|r\|_{2}=1\le 1\) and \(\|A^{\top}r\|_{\infty}\le\gamma\). Strong duality (\cref{rem:p7-strong}) then reads \(\|y^{*}\|_{2}+\gamma\|x^{*}\|_{1}=-b^{\top}r\). Substitute \(\|y^{*}\|_{2}=r^{\top}y^{*}=r^{\top}Ax^{*}-b^{\top}r\) to recover \eqref{eq:p7-comp}. This route uses (a); the directional argument does not.
\end{remark}

\subsubsection{(c) Small columns cannot be used}

The claim does not assume \(Ax^{*}\neq b\). The shortest argument uses only the triangle inequality, and covers both cases.

Suppose \(x_{i}^{*}\neq 0\). For \(0<t<\lvert x_{i}^{*}\rvert\) set \(\sigma=\operatorname{sign}(x_{i}^{*})\) and \(x(t)=x^{*}-t\sigma e_{i}\). Then the \(i\)-th coordinate of \(x(t)\) has absolute value \(\lvert x_{i}^{*}\rvert-t\), and no other coordinate changes, so
\[
\|x(t)\|_{1}=\|x^{*}\|_{1}-t.
\]
The residual becomes \(y^{*}-t\sigma a_{i}\), hence
\[
\|Ax(t)-b\|_{2}
\le
\|y^{*}\|_{2}+t\|a_{i}\|_{2}.
\]
Adding \(\gamma\|x(t)\|_{1}\),
\[
f\bigl(x(t)\bigr)
\le
f(x^{*})+t\bigl(\|a_{i}\|_{2}-\gamma\bigr).
\]
The hypothesis \(\|a_{i}\|_{2}<\gamma\) makes the right-hand side strictly smaller than \(f(x^{*})\) for every small \(t>0\), contradicting optimality of \(x^{*}\). Therefore \(x_{i}^{*}=0\).

If one has already written (b) and is in the case \(y^{*}\neq 0\), a shorter sentence is available: \(x_{i}^{*}\neq 0\) would force \(\lvert a_{i}^{\top}r\rvert=\gamma\) by \eqref{eq:p7-active}, while Cauchy--Schwarz gives \(\lvert a_{i}^{\top}r\rvert\le\|a_{i}\|_{2}\|r\|_{2}=\|a_{i}\|_{2}<\gamma\).

\begin{examnote}[Parts~(b) and~(c) on a script]
\label{examnote:p7-bc}
For (b): expand \(f(x^{*}+te_{i})\) using \eqref{eq:p7-res-diff} and \eqref{eq:p7-l1-diff}, treat \(x_{i}^{*}\neq 0\) and \(x_{i}^{*}=0\) separately, and sum \(u_{i}x_{i}^{*}\). For (c): the triangle argument is one paragraph and does not need \(r\). Do not quote ``soft thresholding'' without writing \(\lvert u_{i}\rvert\le\gamma\).
\end{examnote}

\begin{summary}[Problem~7]
\label{sum:p7}
The Lagrange dual of the residual form is \eqref{eq:p7-dual-prog} (or \eqref{eq:p7-dual-flip}). At a minimizer with nonzero residual, \(A^{\top}r\) lies in the \(\ell_{\infty}\) ball of radius \(\gamma\) and is complementary to \(x^{*}\). A column shorter than \(\gamma\) cannot carry a nonzero coordinate of any minimizer.
\end{summary}


\begin{thebibliography}{99}
\bibitem[GolubVanLoan(1996)]{golub96}
G.~H. Golub and C.~F. Van~Loan.
\newblock \emph{Matrix Computations}, 3rd~ed.
\newblock Johns Hopkins, 1996.
\newblock Official CAM syllabus item~3. Triangular solves and Gaussian elimination: Ch.~3. Matrix inversion via LU: Ch.~3. Cholesky: Ch.~4. Householder QR: Ch.~5.
\bibitem[TrefethenBau(1997)]{trefethen97}
L.~N. Trefethen and D.~Bau, III.
\newblock \emph{Numerical Linear Algebra}.
\newblock SIAM, 1997.
\newblock Official CAM syllabus item~7. SVD: Lectures~4--5. Householder QR: Lectures~10--11. Gaussian elimination and pivoting: Lectures~20--21. Cholesky: Lecture~23. Stability of solving versus inverting: Lecture~21.
\bibitem[ConteDeBoor(2000)]{conte00}
S.~D. Conte and C.~de~Boor.
\newblock \emph{Elementary Numerical Analysis: An Algorithmic Approach}, 3rd~ed.
\newblock McGraw--Hill, 1980; reprint SIAM, 2017.
\newblock Official CAM syllabus item~2. Interpolation remainder. Orthogonal polynomials and Gaussian quadrature.
\bibitem[LeVeque(2007)]{leveque07}
R.~J. LeVeque.
\newblock \emph{Finite Difference Methods for Ordinary and Partial Differential Equations}.
\newblock SIAM, 2007.
\newblock Difference quotients and truncation: Ch.~1. Von Neumann, CFL, Lax equivalence: Chs.~8--10. Du~Fort--Frankel: \S9.3.
\bibitem[Strikwerda(2004)]{strikwerda04}
J.~C. Strikwerda.
\newblock \emph{Finite Difference Schemes and Partial Differential Equations}.
\newblock SIAM, 2004.
\newblock Fourier analysis of difference schemes: Chs.~5--6.
\bibitem[GKO(1995)]{gko95}
B.~Gustafsson, H.-O.~Kreiss, and J.~Oliger.
\newblock \emph{Time Dependent Problems and Difference Methods}.
\newblock Wiley, 1995.
\newblock Official CAM syllabus item~5. Fourier multipliers and the \(1+K\tau\) form of von Neumann stability.
\bibitem[LaxRichtmyer(1956)]{laxrichtmyer56}
P.~D. Lax and R.~D. Richtmyer.
\newblock Survey of the stability of linear finite difference equations.
\newblock \emph{Comm.\ Pure Appl.\ Math.} 9 (1956), 267--293.
\newblock Lax equivalence: a consistent approximation to a linear well-posed IVP converges if and only if it is stable.
\bibitem[BrennerScott(2010)]{brennerscott10}
S.~Brenner and R.~Scott.
\newblock \emph{The Mathematical Theory of Finite Element Methods}.
\newblock Springer, 2010.
\newblock Official CAM syllabus item~8. Weak form, C\'ea. Not the finite-difference language of Problem~4.
\bibitem[BenderOrszag(1999)]{bender99}
C.~M. Bender and S.~A. Orszag.
\newblock \emph{Advanced Mathematical Methods for Scientists and Engineers}.
\newblock Springer, 1999.
\newblock Official CAM syllabus item~1. Regular perturbation of eigenvalues: Ch.~7. Boundary layers: Ch.~9. Multiple scales: Ch.~11.
\bibitem[Nayfeh(1973)]{nayfeh73}
A.~H. Nayfeh.
\newblock \emph{Perturbation Methods}.
\newblock Wiley, 1973.
\newblock Lindstedt--Poincar\'e / strained coordinates: Ch.~4. Multiple scales: Ch.~6.
\bibitem[BoydVandenberghe(2004)]{boyd04}
S.~Boyd and L.~Vandenberghe.
\newblock \emph{Convex Optimization}.
\newblock Cambridge University Press, 2004.
\newblock Dual norms: \S A.1.6. Subgradients: \S3.1.4, \S4.2.3. Lagrange dual, weak duality, Slater and strong duality: \S5.1--5.2.
\end{thebibliography}

\end{document}
